\section{Weak to Strong Cover in Strongly Polynomial Time}\label{sec:WeakToStrongCover}
%\knote{Address the comments from SODA reviews.}\shubhang{done.}
In this section, we 
prove Theorem \ref{thm:WeakToStrongCover:main}. We restate it below. 
%show the following main result:
\thmWeakToStrongCover*

Our proof of \Cref{thm:WeakToStrongCover:main} will require the following definitions. 
\begin{definition}\label{def:(pHw)-tight-sets}
    Let $p:2^V\rightarrow\Z$ be a set function and $(H=(V,E), w)$ be a hypergraph. We will say that the hypergraph $H$ is \emph{non-empty} if $E\neq \emptyset$. 
    %A set $X \subseteq V$ is said to be \emph{tight} w.r.t. the function $p$ and hypergraph $\left(H ,w\right)$ 
    A set $X \subseteq V$ is said to be \emph{$(p, H, w)$-tight} if $b_{(H,w)}(X)=p(X)$. We let $\tightSetFamily{p,H,w}$ denote the family of $(p, H, w)$-tight sets and $\maximalTightSetFamily{p, H, w}$  denote the family of \emph{inclusion-wise maximal} sets of $\tightSetFamily{p,H,w}$. 
    For a hyperedge $e\in E$, we let $\maximalTightSetContainingHyperedge{p, H, w}(e)$ denote an arbitrary set in $\tightSetFamily{p, H, w}$ that contains $e$; if there is no such set, then we will use the convention that $\maximalTightSetContainingHyperedge{p, H, w}(e)$ is undefined.  
    Additionally, we define two operations as follows:
    \begin{enumerate}
        \item For hyperedges $e, f\in E$ and an integer  $\alpha\le \min\{w(e), w(f)\}$, the operation $\merge((H, w), e, f, \alpha)$ returns the hypergraph obtained from $(H, w)$ by decreasing the weight of hyperedges $e$ and $f$ by $\alpha$ and increasing the weight of the hyperedge $e\cup f$ by $\alpha$ (if $\alpha=w(e)$, then we discard $e$; if $\alpha=w(f)$, then we discard $f$; if $e\cup f\not\in E$, then we add the hyperedge $e\cup f$ and set its weight to be $\alpha$). 
        \item For a hyperedge $e\in E$ and an integer $\alpha \le w(e)$, the operation $\reduce((H, w), e, \alpha)$ returns the hypergraph obtained by decreasing the weight of the hyperedge $e$ by $\alpha$ (if $\alpha=w(e)$, then we discard $e$). 
    \end{enumerate}
    %\shubhang{We should either handle the discarding of the 0 weight edges here or in the algorithm. Currently we do it in both places. (we set $E' := \supp(w^R)$ in the algorithm). I am making the edit to handle it here.}
    %For a hyperedge $e \in E$, we let $\maximalTightSetContainingHyperedge{p, H, w}(e)$ denote the \emph{maximal} tight set w.r.t. $(H, w)$ and $p$ containing the hyperedge $e$. We will use the convention that $\maximalTightSetContainingHyperedge{p}(e) \coloneqq  \emptyset$ if there are no tight sets containing $e$.
\end{definition}

We first describe how to prove \Cref{thm:WeakToStrongCover:main} under the assumptions that the input function $p$ and hypergraph $\left(H = (V,E),w\right)$ satisfy the following two technical conditions:
\begin{enumerate}[label=(\alph*), ref=(\alph*)]
    \item \label{cond:WeakToStrong:(a)} %$\maximalTightSetContainingHyperedge{p, H, w}(e) \not = \emptyset$ for every $e \in E$, i.e. 
    every hyperedge $e \in E$ is contained in some $(p, H, w)$-tight set, and
    \item \label{cond:WeakToStrong:(b)} $b_{(H,w)}(u) \not = 0$ for every $u \in V$, i.e. every vertex is contained in some hyperedge.
\end{enumerate}  
Under these assumptions, we will prove \Cref{thm:WeakToStrongCover:main} using the algorithm that we describe in \Cref{sec:WeakToStrongCover:Algorithm} (see \Cref{alg:WeakToStrongCover}). In \Cref{sec:WeakToStrongCover:UncrossingProperties}, we show certain properties of set families that will be used in the description and analysis of our algorithm. In \Cref{sec:WeakToStrongCover:PartialCorrectness}, we show that our algorithm terminates, and the hypergraph returned by the algorithm satisfies properties \ref{thm:WeakToStrongCover:main:(1)} and \ref{thm:WeakToStrongCover:main:(2)} (see \Cref{lem:weak-to-strong:correctness:main:strong-cover} and \Cref{lem:weak-to-strong:correctness:main:hyperedge-merge}). In \Cref{sec:WeakToStrongCover:NumberOfHyperedges}, we show that the hypergraph returned by the algorithm additionally satisfies property \ref{thm:WeakToStrongCover:main:(3)} (see \Cref{lem:WeakToStrongCover:NumberOfHyperedges}). In \Cref{sec:WeakToStrongCover:strongly-poly-runtime:main}, we show that our algorithm can be implemented to run in the claimed strongly polynomial time bound given oracle access to \functionMaximizationOracleStrongCover{p} (see \Cref{lem:WeakToStrong:strongly-polytime:main}).
\Cref{alg:WeakToStrongCover}, \Cref{lem:weak-to-strong:correctness:main:strong-cover}, \Cref{lem:weak-to-strong:correctness:main:hyperedge-merge}, \Cref{lem:WeakToStrongCover:NumberOfHyperedges}, and \Cref{lem:WeakToStrong:strongly-polytime:main} together complete the proof of \Cref{thm:WeakToStrongCover:main} under the assumptions that conditions \ref{cond:WeakToStrong:(a)} and \ref{cond:WeakToStrong:(b)} hold.

We now describe how to circumvent assumptions \ref{cond:WeakToStrong:(a)} and \ref{cond:WeakToStrong:(b)} on the input function $p$ and hypergraph $(H,w)$. Let $e \in E$ be a hyperedge not contained in any $(p, H, w)$-tight set. Let $\alpha_e^R \coloneqq  \min\{w(e), \beta^R(e)\}$, where $$\beta^R(e) \coloneqq  \min\{b_{(H,w)}(X) - p(X) : e \subseteq X \subseteq V\}.$$
We construct the hypergraph $(H_0^e:=(V, E_0^e:=\{e\}), w_0^e:E_0\rightarrow \{\alpha^R_e\})$. Now, consider the function $p' \coloneqq  p - d_{(H_0^e, w_0^e)}$ and the  hypergraph $(H', w'):=\reduce((H, w), e, \alpha^R_e)$; we term this operation as the reduce operation with respect to the hyperedge $e$. 
%(jumping ahead, we note that this operation is the \emph{reduce} operation which we describe in the next \Cref{sec:WeakToStrongCover:Algorithm}).
Iteratively performing this reduce operation with respect to every hyperedge $e \in E$ results in a function $p_1$ and a hypergraph $(H_1 = (V, E_1), w_1)$ such that $b_{(H_1, w_1)}(X)\ge p_1(X)$ for every $X\subseteq V$ and also satisfies condition \ref{cond:WeakToStrong:(a)}. Furthermore, suppose that $\zeros \coloneqq  \{u \in V : b_{H_1, w_1}(u) = 0\}$. Then, the function $p_2 \coloneqq  \functionContract{p_1}{\zeros}$ and hypergraph $(H_2 = (V - \zeros, E_1), w_1)$ is an instance with $b_{(H_2, w_1)}(X)\ge p_2(X)$ for every $X\subseteq V$ and also satisfies conditions \ref{cond:WeakToStrong:(a)} and \ref{cond:WeakToStrong:(b)}. Finally, we observe that if $(H^*, w^*)$ is a hypergraph that strongly covers the function $p_2$, then the hypergraph $(H^*+\sum_{e \in E}H_0^e, w^*+w_0^e)$ strongly covers the original function $p$. Moreover,  if $(H^*, w^*)$ is obtained by merging hyperedges of $(H_2, w_1)$, then the  hypergraph $(H^*+\sum_{e \in E}H_0^e, w^*+w_0^e)$ is obtained by merging hyperedges of $(H, w)$. Thus, our final algorithm is to first obtain the instance $(H_2, w_1)$ and an oracle for \functionMaximizationOracleStrongCover{p_2} using \functionMaximizationOracleStrongCover{p} such that the instance $(p_2, H_2, w_1)$ is such that $(H_2, w_1)$ weakly covers $p_2$ and satisfies the technical conditions \ref{cond:WeakToStrong:(a)} and \ref{cond:WeakToStrong:(b)}, obtain a hypergraph $(H^*, w^*)$ by the procedure described in the previous paragraph, and finally return $(H^*+\sum_{e \in E}H_0^e, w^*+w_0^e)$. 
%\kristof{We should define this operation, that is, the `addition' of edge-weighted hypergraphs in the preliminaries.}\shubhang{todo.}

\subsection{Weak to Strong Cover Algorithm}\label{sec:WeakToStrongCover:Algorithm}
Our algorithm is recursive and takes two inputs: (1) a hypergraph $(H = (V, E), w)$, and (2) a symmetric skew-supermodular function $p:2^V \rightarrow\Z$. We note that in contrast to the statement of \Cref{thm:WeakToStrongCover:main}, the inputs $(H,w)$ and $p$ to the algorithm are required to satisfy the two technical conditions 
\ref{cond:WeakToStrong:(a)} and \ref{cond:WeakToStrong:(b)} specified in the previous section. The algorithm returns a hypergraph $\left(H^*= \left(V, E^*\right), w^*\right)$. 

We now give an informal description of the algorithm (refer to \Cref{alg:WeakToStrongCover} for a formal description). If $E=\emptyset$, then the algorithm returns the empty hypergraph. Otherwise, the algorithm chooses two hyperedges $e, f \in E$  such that $\maximalTightSetContainingHyperedge{p, H, w}(e)\neq \maximalTightSetContainingHyperedge{p, H, w}(f)$ ($\maximalTightSetContainingHyperedge{p, H, w}(e)$ is well-defined for every hyperedge $e\in E$ owing to condition \ref{cond:WeakToStrong:(a)}; in the next section, we show that $\maximalTightSetContainingHyperedge{p, H, w}(e)$ is unique for every hyperedge $e\in E$ under condition \ref{cond:WeakToStrong:(a)} and also that if $(H, w)$ is not a strong cover of $p$, then there exists a pair of hyperedges $e, f\in E$ such that $\maximalTightSetContainingHyperedge{p, H, w}(e)\neq \maximalTightSetContainingHyperedge{p, H, w}(f)$). 
Next, 
%the algorithm performs \merge and \reduce operations (see \Cref{def:(pHw)-tight-sets}). In particular, 
the algorithm computes the value 
 $$\beta^M \coloneqq  \min\left\{b_{(H,w)}(X) - p(X) : X \subseteq V \text{ and } e \cap X, f \cap X \not = \emptyset\right\},$$ sets $\alpha^M \coloneqq  \min\{\beta^M, w(e), w(f)\}$, and defines the hypergraph $(H^M, w^M) := \merge((H, w), e, f, \alpha^M)$.
%The algorithm then computes an intermediate weight function $w^M \coloneqq  w - \alpha^M\chi_e -\alpha^M\chi_f + \alpha^M\chi_{e\cup f}$, where $\chi_e, \chi_f, \chi_{e\cup f}\in \{0, 1\}^{E\cup\{e\cup f\}}$ are the indicator vectors of hyperedges $e, f, e\cup f$. We will refer to these two steps as \emph{merge} steps. 
Next, the algorithm computes the value $$\beta^R \coloneqq  \min\left\{b_{(H^M, w^M)}\left(X\right) - p\left(X\right) : e \cup f \subseteq X \subseteq V\right\},$$ 
    sets $\alpha^R \coloneqq  \min\{\beta^R, w^M(e\cup f)\}$, and defines the hypergraph $(H^R, w^R) := \reduce((H^M, w^M), e \cup f, \alpha^R)$.
%The algorithm then computes the new weight function $w^R \coloneqq  w^M - \alpha^R\chi_{e\cup f}$. We will refer to these two steps as \emph{reduce} steps.
It computes hypergraph $(H_0=(V, E_0:=\{e\cup f\}), w_0: E_0\rightarrow \{\alpha^R\})$ and defines $p^R:=p-d_{(H_0, w_0)}$. 
Next, the algorithm computes the set $\zeros \coloneqq  \{u \in V : b_{(H^R, w^R)}(u) = 0\}$, and recursively calls itself with the inputs 
    $(H' = (V':=V-\zeros, E':=E^R),w':=w^R) $ and $p' \coloneqq  \functionContract{p^R}{\zeros}$ to obtain a hypergraph $\left(H_0^*, w_0^*\right)$. Finally, the algorithm obtains the hypergraph $(G, c)$ from $(H_0^*, w_0^*)$ by adding the set $\zeros$ of vertices and returns the hypergraph $(G+H_0, c + w_0)$.
    %, adds the hyperedge $e\cup f$ with (additional) weight $\alpha^R$, and returns the resulting hypergraph $(G, c) := $. 
    %We formally describe our algorithm in \Cref{alg:WeakToStrongCover}.


\begin{algorithm}[!h]
\caption{Weak to Strong Cover}\label{alg:WeakToStrongCover}
$\mathbf{\textsc{Algorithm}}\left(H = (V, E), w, p\right):$
\begin{algorithmic}[1]
    \If{$E = \emptyset$} \textsc{Base Case:}
    \State{\Return{$(H, w)$.}}
    \Else:
    \State{Pick a pair of hyperedges $e, f \in E $ s.t. $\maximalTightSetContainingHyperedge{p, H, w}(e) \not = \maximalTightSetContainingHyperedge{p, H, w}(f)$}\label{alg:WeakToStrongCover:(4)}
    \State{$\alpha^M \coloneqq  \min\begin{cases}
    \beta^M \coloneqq  \min\left\{b_{(H,w)}(X) - p(X) : X \subseteq V \text{ and } e \cap X, f \cap X \not = \emptyset\right\} & \\
    \min\left\{w(e), w(f)\right\}&
    \end{cases}
    $}\label{alg:WeakToStrongCover:(5)}
    % \vspace*{0.5mm}
    %\State{$\alpha^M \coloneqq  \left\{w(e), w(f), \min\left\{b_{(H,w)}(X) - p(X) : e\cup f \subseteq X \subseteq V\right\}\right\}$}
    \State{$(H^M, w^M):=\merge\left(\left(H, w\right), e, f, \alpha^M\right)$}\label{alg:WeakToStrongCover:(6)}
    %\State{Let $(H^M, w^M)$ be the hypergraph obtained from $(H, w)$ by reducing the weight of hyperedges $e$ and $f$ by $\alpha^M$ and increasing the weight of the hyperedge $e\cup f$ by $\alpha^M$ (discard hyperedge $e$ and/or $f$ if the weight of that hyperedge has decreased to zero;  introduce a new hyperedge $e\cup f$ with weight $\alpha^M$ if $e\cup f\not\in E$)}\label{alg:WeakToStrongCover:(6)}
    %\State{$w^M \coloneqq  w - \alpha^M\chi_e - \alpha^M\chi_f + \alpha^M\chi_{e\cup f}$}\label{alg:WeakToStrongCover:(6)}
    %\vspace*{0.5mm}
    \State{$\alpha^R \coloneqq  \min \begin{cases}
        \beta^R \coloneqq  \min\left\{b_{(H, w^M)}\left(X\right) - p\left(X\right) : e \cup f \subseteq X \subseteq V\right\} &\\
        w^M(e \cup f) &
    \end{cases}$}\label{alg:WeakToStrongCover:(7)}
    %\State{$\alpha^R \coloneqq  \min\left\{w^M(e \cup f), \min\left\{b_{(H^M, w^M)}(X) - p(X) : e\cup f \subseteq X \subseteq V\right\}\right\}$}
    %\vspace*{0.5mm}
    \State{$(H^R, w^R) \coloneqq \reduce\left(\left(H^M, w^M\right), e\cup f, \alpha^R\right)$}\label{alg:WeakToStrongCover:(8)}
    %\vspace*{0.5mm}
    \State{Construct $\left(H_0:=\left(V, E_0:=\left\{e\cup f\right\}\right), w_0: E_0\rightarrow \left\{\alpha^R\right\}\right)$}\label{alg:WeakToStrongCover:(9)}
    \State{$p^R\coloneqq p - d_{(H_0, w_0)}$}\label{alg:WeakToStrongCover:(10)}
    %\State{$w^R \coloneqq  w^M - \alpha^R\chi_{e\cup f}$ and $p^R\coloneqq p - \alpha^R d_{(V, \{e\cup f\})}$}\label{alg:WeakToStrongCover:(8)}
    %\vspace*{0.5mm}
    
    \State{$\zeros \coloneqq  \left\{v \in V : b_{(H^R, w^R)}(v) = 0\right\}$}\label{alg:WeakToStrongCover:(11)}
    % \vspace*{1mm}
    
    \State{$V' \coloneqq  V - \zeros$}\label{alg:WeakToStrongCover:(12)}
    \State{$p' \coloneqq  \functionContract{p^R}{\zeros}$}\label{alg:WeakToStrongCover:(13)}
    
    \State{$(H_0^*, w_0^*):=\textsc{Algorithm}\left(H' = \left(V', E' := E^R\right), w' := w^R , p'\right)$}\label{alg:WeakToStrongCover:(14)}
    
    \State{Obtain $(G, c)$ from $(H_0^*, w_0^*)$ by adding vertices \zeros}\label{alg:WeakToStrongCover:(15)}
    
    \State{\Return{$(G+H_0, c+w_0 )$}}\label{alg:WeakToStrongCover:(16)}
    \EndIf
\end{algorithmic}
\end{algorithm}

\subsection{Properties of Tight Sets}\label{sec:WeakToStrongCover:UncrossingProperties}
%\knote{Can this section be folded into Section 4 and the entire section be moved before Section 3?}
In this section, we show properties of the set families $\tightSetFamily{p, H, w}$ and $\maximalTightSetFamily{p, H, w}$ which will be useful to show correctness and convergence of \Cref{alg:WeakToStrongCover} in subsequent sections.
%We recall that  a set $X\subseteq V$ is said to be \emph{tight} w.r.t. a function $p:2^V\rightarrow\Z$ and a hypergraph $(H = (V, E), w)$ if $p(X) = b_{(H,w)}(X)$. We let $\tightSetFamily{p,H,w}$ denote the family of tight sets and $\maximalTightSetFamily{p, H, w}$ denote the family of \emph{inclusion-wise maximal} tight sets w.r.t. the function $p$ and hypergraph $(H,w)$. We also recall that for a hyperedge $e \in E$, the set $\maximalTightSetContainingHyperedge{p, H, w}(e)$ denotes the \emph{maximal} set in $\tightSetFamily{p, H, w}$ containing $e$. 
%\knote{Include a claim that $\maximalTightSetContainingHyperedge{p, H, w}(e)$ is unique for every hyperedge $e$.} 
%We use the convention that $\maximalTightSetContainingHyperedge{p, H, w}(e) \coloneqq  \emptyset$ if no such set exists. 
In particular, \Cref{lem:weak-to-strong:uncrossing:uncrossing-lemma} below shows an important property of the family of tight sets and \Cref{cor:weak-to-strong:uncrossing:maximal-tight-set-containing-hyperedge-unique} shows that for every hyperedge $e\in E$, there exists a unique maximal $(p, H, w)$-tight set containing $e$ under condition \ref{cond:WeakToStrong:(a)}. Furthermore, \Cref{lem:weak-to-strong:uncrossing:maximal-tight-set-family-disjoint} shows that the family $\maximalTightSetFamily{p, H, w}$ is disjoint under condition \ref{cond:WeakToStrong:(b)}, and consequently, \Cref{cor:weak-to-strong:uncrossing:exist-two-hyperedges-contained-in-disjoint-maximal-tight-sets} implies  that \Cref{alg:WeakToStrongCover} \Cref{alg:WeakToStrongCover:(4)} is well-defined.

\begin{lemma}\label{lem:weak-to-strong:uncrossing:uncrossing-lemma}
Let $p:2^V\rightarrow\Z$ be a skew-supermodular function and let $\left(H = (V, E), w:E \rightarrow\Z_+\right)$ be a hypergraph such that $b_{(H,w)}(Z) \geq p(Z)$ for every $Z \subseteq V$. Let $X, Y \in \tightSetFamily{p, H, w}$ be sets such that $X\cap Y \not = \emptyset$.
If there exists a hyperedge $e \in E$ such that (1) $e$ intersects the set $X\cap Y$ and 
(2) $e$ intersects at most one of the sets $X - Y$ and $Y - X$, then $X\cup Y, X \cap Y \in \tightSetFamily{p, H, w}$.
\end{lemma}
\begin{proof}
We note that the claim holds if $X\subseteq Y$ or $Y \subseteq X$. Suppose that $X - Y, Y - X \not= \emptyset$. First, suppose that the function $p$ satisfies $p(X) + p(Y) \leq p(X - Y) + p(Y - X)$. Then, the following gives us a contradiction:
\begin{align*}
    b_{(H,w)}(X) + b_{(H,w)}(Y) 
    &= p(X) + p(Y) &\\
    & \leq p(X - Y) + p(Y - X) & \\
    & \leq b_{(H,w)}(X - Y) + b_{(H,w)}(Y - X) & \\
    & < b_{(H,w)}(X) + b_{(H,w)}(Y).
\end{align*}
Here, the first equality is because $X, Y \in \tightSetFamily{p, H, w}$. The second inequality is because $b_{(H,w)}(Z) \geq p(Z)$ for every $Z\subseteq V$. The final inequality is because the hyperedge $e$ intersects at most one of the sets $X - Y$ and $Y - X$.
Thus, the function $p$ satisfies $p(X) + p(Y) \leq p(X\cap Y) + p(X \cup Y)$. Consequently, we have that
\begin{align*}
    b_{(H,w)}(X) + b_{(H,w)}(Y) 
    &= p(X) + p(Y) &\\
    & \leq p(X \cup Y) + p(X \cap Y) &\\
    & \leq b_{(H,w)}(X\cup Y) + b_{(H,w)}(X\cap Y) &\\
    & \leq b_{(H,w)}(X) + b_{(H,w)}(Y).&
\end{align*}
Here, the first equality is because $X, Y \in \tightSetFamily{p, H, w}$. The second inequality is because $b_{(H,w)}(Z) \geq p(Z)$ for every $Z\subseteq V$. The final inequality is because the coverage function $b_{(H, w)}$ is submodular. Thus, all inequalities in the above sequence are in fact equalities, and consequently we have that $X\cup Y, X\cap Y \in \tightSetFamily{p, H, w}$.
\end{proof}

\begin{corollary}\label{cor:weak-to-strong:uncrossing:maximal-tight-set-containing-hyperedge-unique}
Let $p:2^V\rightarrow\Z$ be a skew-supermodular function and let $\left(H = (V, E), w:E \rightarrow\Z_+\right)$ be a hypergraph such that  
every hyperedge $e \in E$ is contained in some $(p, H, w)$-tight set
%condition \ref{cond:WeakToStrong:(a)} such that 
and $b_{(H,w)}(Z) \geq p(Z)$ for every $Z \subseteq V$. 
%and every hyperedge $e \in E$ is contained in some set from the family $\tightSetFamily{p, H, w}$. 
Then, for every hyperedge $e \in E$, there is a unique set from the family $\maximalTightSetFamily{p, H, w}$ containing the hyperedge $e$.
\end{corollary}
\begin{proof}
    Let $e \in E$ be a hyperedge. We note that there exists a $(p, H, w)$-tight set containing the hyperedge $e$. Consequently, there exists a set in the family \maximalTightSetFamily{p, H, w} containing the hyperedge $e$. By way of contradiction, suppose that $X, Y \in \maximalTightSetFamily{p, H, w}$ are distinct sets containing the hyperedge $e$. We note that $X - Y, Y - X \not = \emptyset$ by the maximality of the sets $X, Y \in \maximalTightSetFamily{p, H, w}$.
    Also, $X\cap Y\neq \emptyset$ since $e$ is contained in $X\cap Y$. Moreover, $e$ is disjoint from both $X-Y$ and $Y-X$. Hence, by \Cref{lem:weak-to-strong:uncrossing:uncrossing-lemma}, we have that $X\cup Y \in \tightSetFamily{p, H, w}$, contradicting maximality of the sets $X, Y \in \maximalTightSetFamily{p, H, w}$.
\end{proof}

\iffalse
By \Cref{cor:weak-to-strong:uncrossing:maximal-tight-set-containing-hyperedge-unique}, $\maximalTightSetContainingHyperedge{p, H, w}(e)$ denotes the  unique maximal $(p, H, w)$-tight set containing a hyperedge $e\in E$ when condition \ref{cond:WeakToStrong:(a)} holds.
\begin{definition}
Let $p:2^V\rightarrow\Z$ be a skew-supermodular function and let $\left(H = (V, E), w:E \rightarrow\Z_+\right)$ be a hypergraph such that $b_{(H,w)}(Z) \geq p(Z)$ for every $Z \subseteq V$ and every hyperedge $e \in E$ is contained in some $(p, H, w)$-tight set. For a hyperedge $e \in E$, we let $\maximalTightSetContainingHyperedge{p, H, w}(e) \in \maximalTightSetFamily{p, H, w}$ denote the \emph{unique maximal} $(p, H, w)$-tight containing the hyperedge $e$. 
    %We will use the convention that $\maximalTightSetContainingHyperedge{p, H, w}(e) \coloneqq  \emptyset$ if there are no tight sets containing $e$.
\end{definition}
\fi
The next lemma shows that for a \emph{skew-supermodular} function $p$ and hypergraph $(H, w)$ satisfying conditions  \ref{cond:WeakToStrong:(a)} and \ref{cond:WeakToStrong:(b)}, the family $\maximalTightSetFamily{p, H, w}$ is disjoint. The corollary following the lemma says that if the function $p$ is also \emph{symmetric} and the hypergraph $H$ is non-empty, then there exist two hyperedges that are each contained in distinct maximal $(p, H, w)$-tight sets w.r.t. $p$ and $(H, w)$.  We note that this makes \Cref{alg:WeakToStrongCover} \Cref{alg:WeakToStrongCover:(4)} well-defined.

\begin{lemma}\label{lem:weak-to-strong:uncrossing:maximal-tight-set-family-disjoint}
    Let $p:2^V\rightarrow\Z$ be a skew-supermodular function and let $\left(H = (V, E), w:E \rightarrow\Z_+\right)$ be a hypergraph 
    %satisfying conditions \ref{cond:WeakToStrong:(a)} and \ref{cond:WeakToStrong:(b)} 
    such that every hyperedge $e \in E$ is contained in some $(p, H, w)$-tight set, $b_{(H,w)}(u) \not = 0$ for every $u \in V$, and $b_{(H,w)}(X) \geq p(X)$ for every $X \subseteq V$. 
    %Furthermore, suppose that $\maximalTightSetContainingHyperedge{p, H, w}(e) \not = \emptyset$ for every $e \in E$ and $b_{(H,w)}(u) > 0$ for every $u \in V$. 
    Then, the family $\maximalTightSetFamily{p, H, w}$ is a disjoint family.
\end{lemma}
\begin{proof}
By way of contradiction, let $X, Y \in \maximalTightSetFamily{p, H, w}$ be distinct sets such that $X\cap Y \not = \emptyset$. By maximality of the sets $X, Y$, we also have that $X - Y, Y - X \not = \emptyset$. Let $u \in X\cap Y$. Since $b_{(H,w)}(u) > 0$, there exists a hyperedge $e \in E$ such that $u \in e$. We note that if either $e\cap (X - Y) = \emptyset$ or $e\cap (Y - X) = \emptyset$, then the set $X \cup Y \in \tightSetFamily{p, H, w}$ by \Cref{lem:weak-to-strong:uncrossing:uncrossing-lemma}, contradicting maximality of $X, Y \in \maximalTightSetFamily{p, H, w}$. Thus, we have that $e \cap (X - Y), e \cap (Y - X) \not= \emptyset$. Furthermore, since the hyperedge $e$ is contained in some $(p, H, w)$-tight set, we have that  $\maximalTightSetContainingHyperedge{p, H, w}(e) \not = \emptyset$. Consequently, we have that $\maximalTightSetContainingHyperedge{p, H, w}(e)\cap (X - Y), \maximalTightSetContainingHyperedge{p, H, w}(e) \cap (Y - X) \not = \emptyset$. However, $e \subseteq \maximalTightSetContainingHyperedge{p, H, w}(e)$ and thus, $e \cap (X - \maximalTightSetContainingHyperedge{p, H, w}(e)) = \emptyset$. By \Cref{lem:weak-to-strong:uncrossing:uncrossing-lemma}, we have that the set $X \cup \maximalTightSetContainingHyperedge{p, H, w}(e) \in \tightSetFamily{p, H, w}$, contradicting maximality of the set $X \in \maximalTightSetFamily{p, H, w}$.
\end{proof}

We recall that a hypergraph is non-empty if it contains at least one hyperedge. 
\begin{corollary}[\hspace{-1sp}\cite{Bernath-Kiraly}]
\label{cor:weak-to-strong:uncrossing:exist-two-hyperedges-contained-in-disjoint-maximal-tight-sets}
Let $p:2^V\rightarrow\Z$ be a symmetric skew-supermodular function and $\left(H = (V, E), w:E \rightarrow\Z_+\right)$ be a non-empty hypergraph 
%\knote{Non-emptiness is not needed since $b_{(H, w)}(u)\neq 0$ for every $u\in V$?}\shubhang{$b_{(H, w)}(u)\neq 0$ for every $u\in V$ still allows $V = \emptyset$.} \knote{Ensure that the use of the corollary guarantees this hypothesis.} 
%satisfying conditions \ref{cond:WeakToStrong:(a)} and \ref{cond:WeakToStrong:(b)} 
such that every hyperedge $e \in E$ is contained in some $(p, H, w)$-tight set, $b_{(H,w)}(u) \not = 0$ for every $u \in V$, and $b_{(H,w)}(X) \geq p(X)$ for every $X \subseteq V$. 
%Let $p:2^V\rightarrow\Z$ be a symmetric skew-supermodular function. Furthermore, let $\left(H = (V, E), w:E \rightarrow\Z\right)$ be a non-empty hypergraph such that $b_{(H,w)}(X) \geq p(X)$ for every $X \subseteq V$ such that $\maximalTightSetContainingHyperedge{p, H, w}(e) \not = \emptyset$ for every hyperedge $e \in E$ and $b_{(H,w)}(u) > 0$ for every vertex $u \in V$. 
Then, there exist distinct hyperedges $e,f \in E$ such that $\maximalTightSetContainingHyperedge{p, H, w}(e) \cap \maximalTightSetContainingHyperedge{p, H, w}(f) = \emptyset$.
\end{corollary}
\begin{proof}
Let $e \in E$ be a hyperedge.
We will show that that there exists a hyperedge $f \in E$ such that $f \subseteq V - \maximalTightSetContainingHyperedge{p, H, w}(e)$. The claim will then follow by \Cref{lem:weak-to-strong:uncrossing:maximal-tight-set-family-disjoint} since $\maximalTightSetContainingHyperedge{p, H, w}(f) \not = \emptyset$ (because the hyperedge $f$ is contained in some $(p, H, w)$-tight set) and $\maximalTightSetContainingHyperedge{p, H, w}(e), \maximalTightSetContainingHyperedge{p, H , w}(f) \in \maximalTightSetFamily{p, H, w}$.
%$, the claim follows from \Cref{lem:weak-to-strong:uncrossing:maximal-tight-set-family-disjoint}. 
By way of contradiction, suppose that there is no hyperedge contained in the set $V - \maximalTightSetContainingHyperedge{p, H, w}(e)$. Then, we have that $p(V - \maximalTightSetContainingHyperedge{p, H, w}(e)) = p(\maximalTightSetContainingHyperedge{p, H, w}(e)) = b_{(H,w)}(\maximalTightSetContainingHyperedge{p, H, w}(e)) > b_{(H,w)}(V - \maximalTightSetContainingHyperedge{p, H, w}(e))$, a contradiction. Here, the first equality is because the function $p$ is symmetric and the final inequality is because every hyperedge that intersects the set $V - \maximalTightSetContainingHyperedge{p, H, w}(e)$ also intersects the set $\maximalTightSetContainingHyperedge{p, H, w}(e)$ by our contradiction assumption.
\end{proof}


\subsection{Termination and Partial Correctness of the Algorithm}
%\Cref{alg:WeakToStrongCover}}
\label{sec:WeakToStrongCover:PartialCorrectness}
In this section, we show that every recursive call of \Cref{alg:WeakToStrongCover} is well-defined. 
%when the inputs to the base recursive call satisfy the aforementioned properties. 
Furthermore, we show that \Cref{alg:WeakToStrongCover} terminates and that the hypergraph returned by the algorithm satisfies properties \ref{thm:WeakToStrongCover:main:(1)} and \ref{thm:WeakToStrongCover:main:(2)} of \Cref{thm:WeakToStrongCover:main}. The main results of this section are \Cref{lem:weak-to-strong:correctness:main:strong-cover} and \Cref{lem:weak-to-strong:correctness:main:hyperedge-merge}. We prove certain additional lemmas in this section which will be useful in subsequent sections. 
%We will use the convention that $\Tilde{e}$ refers to an arbitrary hyperedge in $E$, where $E$ is the hyperedge set of the hypergraph $H$ input to the current recursive call and use $e, f$ to denote the hyperedges picked by \Cref{alg:WeakToStrongCover} \Cref{alg:WeakToStrongCover:(4)} in the current recursive call. \knote{The last sentence can be deleted.}
 
    Suppose that the input to \Cref{alg:WeakToStrongCover} is a non-empty hypergraph $(H = (V, E), w)$ and a symmetric skew-supermodular function $p:2^V\rightarrow \Z$ 
%satisfying conditions \ref{cond:WeakToStrong:(a)} and \ref{cond:WeakToStrong:(b)} 
such that every hyperedge $e \in E$ is contained in some $(p, H, w)$-tight set, $b_{(H,w)}(u) \not = 0$ for every $u \in V$, and $b_{(H,w)}(X) \geq p(X)$ for every $X \subseteq V$. 
  Then, by \Cref{cor:weak-to-strong:uncrossing:exist-two-hyperedges-contained-in-disjoint-maximal-tight-sets}, there exist hyperedges $e , f \in E$ such that $\maximalTightSetContainingHyperedge{p, H, w}(e) \not = \maximalTightSetContainingHyperedge{p, H, w}(f)$. In particular, \Cref{alg:WeakToStrongCover} \Cref{alg:WeakToStrongCover:(4)} is well-defined. 
  %In this section, we show that every recursive call of \Cref{alg:WeakToStrongCover} is well-defined when the inputs to the base recursive call satisfy the aforementioned properties. Furthermore, we show that \Cref{alg:WeakToStrongCover} terminates and that the hypergraph returned by the algorithm satisfies properties (1) and (2) of \Cref{thm:WeakToStrongCover:main}. The main results of this section are \Cref{lem:weak-to-strong:correctness:main:strong-cover} and \Cref{lem:weak-to-strong:correctness:main:hyperedge-merge}. 
  The next two lemmas establish several useful properties of the intermediate quantities used in the algorithm.

\begin{lemma}\label{lem:WeakToStrong:properties-of-alpha^M-and-alpha^R}
    %Consider a recursive call of \Cref{alg:WeakToStrongCover}. 
    Suppose that the input to \Cref{alg:WeakToStrongCover} is a non-empty hypergraph $(H = (V, E), w)$ and a symmetric skew-supermodular function $p:2^V\rightarrow \Z$ 
such that every hyperedge $\Tilde{e} \in E$ is contained in some $(p, H, w)$-tight set, $b_{(H,w)}(u) \not = 0$ for every $u \in V$, and $b_{(H,w)}(X) \geq p(X)$ for every $X \subseteq V$. 
    %Furthermore, suppose that $\maximalTightSetContainingHyperedge{p, H, w}(\Tilde{e}) \not = \emptyset$ for every hyperedge $\Tilde{e} \in E$ and $b_{(H,w)}(u) > 0$ for every vertex $u \in V$. 
    Let $e, f, \alpha^M, (H^M, w^M), \alpha^R$, $(H^R, w^R),(H_0, w_0), p^R$ and $(H' = (V', E'), w')$ be as defined by \Cref{alg:WeakToStrongCover}. 
    %Let $(H' = (V', E'), w^R)$ be the hypergraph and $p':2^{V'}\rightarrow\Z$ be the function computed by \Cref{alg:WeakToStrongCover}. 
    Then, we have the following:
    \begin{enumerate}[label=(\arabic*)]
        \item $\alpha^M \in \Z_+$,
        \item 
        $E^M=E-\{e\}\cdot \indicator_{\alpha^M= w(e)}- \{f\}\cdot \indicator_{\alpha^M = w(f)} + \{e\cup f\}$ with the weights of the hyperedges in $E^M$ being  
        \[
        w^M(\Tilde{e})=
        \begin{cases}
            w(\Tilde{e}) &\text { if } \Tilde{e}\in E-\{e, f, e\cup f\},\\
            w(e\cup f) + \alpha^M &\text { if } \Tilde{e}=e\cup f \text{ and } e\cup f \in E, \\
            \alpha^M &\text { if } \Tilde{e}=e\cup f \text{ and } e\cup f \not \in E, \\
            w(e) - \alpha^M &\text { if } \Tilde{e}=e \text{ and } \alpha^M<w(e), \\
            w(f) - \alpha^M &\text { if } \Tilde{e}=f \text{ and } \alpha^M<w(f).
        \end{cases}
        \]
        %$w^M = w - w_0^M + w_1^M$, where $w_0^M:\{e, f\}\rightarrow\{\alpha^M\}$ and $w_1^M:\{e\cup f\} \rightarrow\{\alpha^M\}$, 
        %\knote{I am not a fan of this notation. Instead, can we say: $E^M=E-\{e\}\cdot \indicator_{\alpha^M= w(e)}- \{f\}\cdot \indicator_{\alpha^M = w(f)} + \{e\cup f\}$ with the weights of the hyperedges in $E^M$ being  
        % \[
        % w^M(\Tilde{e})=
        % \begin{cases}
        %     w(\Tilde{e}) &\text { if } \Tilde{e}\in E-\{e, f, e\cup f\},\\
        %     w(e\cup f) + \alpha^M &\text { if } \Tilde{e}=e\cup f \text{ and } e\cup f \in E, \\
        %     \alpha^M &\text { if } \Tilde{e}=e\cup f \text{ and } e\cup f \not \in E, \\
        %     w(e) - \alpha^M &\text { if } \Tilde{e}=e \text{ and } \alpha^M<w(e), \\
        %     w(f) - \alpha^M &\text { if } \Tilde{e}=f \text{ and } \alpha^M<w(f).
        % \end{cases}
        % \]
        % }
        \item $b_{(H^M, w^M)}(X) \geq p(X)$ for every $X \subseteq V$,
        \item $\alpha^R \in \Z_{\geq 0}$, 
        \item 
        $E' = E^R = E^M - \{e\cup f\}\cdot \indicator_{\alpha^R=w^M(e\cup f)}$ with the weights of the hyperedges in $E' = E^R$ being 
        \[
        w'(\Tilde{e}) = w^R(\Tilde{e})=
        \begin{cases}
            w^M(\Tilde{e}) &\text { if } \Tilde{e}\in E^M-\{e\cup f\},\\
            w^M(e\cup f) - \alpha^R &\text { if } \Tilde{e}=e\cup f \text{ and } \alpha^R< w^M(e\cup f), \text{ and}  \\
        \end{cases}
        \]
        %$w^R = w^M - w_0$, \knote{Again, not a fan of the notation. Instead, can we say: 
        % $E^R = E^M - \{e\cup f\}\cdot \indicator_{\alpha^R<w^M(e\cup f)}$ with the weights of the hyperedges in $E^R$ being 
        % \[
        % w^R(\Tilde{e})=
        % \begin{cases}
        %     w^M(\Tilde{e}) &\text { if } \Tilde{e}\in E^M-\{e\cup f\},\\
        %     w^M(e\cup f) - \alpha^R &\text { if } \Tilde{e}=e\cup f \text{ and } \alpha^R< w^M(e\cup f), \\
        % \end{cases}
        % \]
        % }
        \item $b_{(H^R, w^R)}(X) \geq p^R(X)$ for every $X \subseteq V$. 
        %\item $\alpha^R = 0$ if and only if $\alpha^M = \beta^M$ ( i.e., $\alpha^R > 0$ if and only if $\alpha^M < \beta^M$) , and
        % \item $E' = E - \{e\}\cdot \indicator_{\alpha^M = w(e)} - \{f\}\cdot \indicator_{\alpha^M = w(f)} + \{e\cup f\}\cdot \indicator_{\alpha^R < w^M(e\cup f)}$.
        %- \{e, f\} \subseteq E' \subseteq E \cup \{e \cup f\}$,
        % \item if $\alpha^R = \beta^R < w^M(e\cup f)$, then $e \cup f \in E'$,
        % \item if $\alpha^R < \beta^R$, then $E' \subseteq E$.
    \end{enumerate}
\end{lemma}
\begin{proof} We prove each claim separately below.
\begin{enumerate}
    \item We recall that the weight function $w:E\rightarrow\Z_+$ is a positive integer-valued function, and the function $p$ is an integer-valued function. Thus, it suffices to show that $\beta^M > 0$. We note that $\beta^M \geq 0$ since $b_{(H,w)}(X) \ge p(X)$ for every $X \subseteq V$. By way of contradiction, suppose that $\beta^M = 0$. Let $X \subseteq V$ witness $\beta^M$, i.e. $e \cap X, f\cap X \not = \emptyset$ and $b_{(H,w)}(X) -p(X)= 0$.
    The latter property implies that $X \in \tightSetFamily{p, H, w}$ and thus, there exists $Y \in \maximalTightSetFamily{p, H, w}$ such that $X \subseteq Y$. 
    We recall that by the algorithm's choice of hyperedges $e, f \in E$, we have that $\maximalTightSetContainingHyperedge{p, H, w}(e) \not = \maximalTightSetContainingHyperedge{p, H, w}(f)$. Since, $\maximalTightSetContainingHyperedge{p, H, w}(e), \maximalTightSetContainingHyperedge{p, H, w}(f) \in \maximalTightSetFamily{p, H, w}$, we have that $\maximalTightSetContainingHyperedge{p, H, w}(e) \cap \maximalTightSetContainingHyperedge{p, H, w}(f) = \emptyset$ by \Cref{lem:weak-to-strong:uncrossing:maximal-tight-set-family-disjoint}. In particular, we have that $\maximalTightSetContainingHyperedge{p, H, w}(e) \not = Y$ since
    $f \subseteq V - \maximalTightSetContainingHyperedge{p, H, w}(e)$ but $f \cap Y \not = \emptyset$.
    Thus, $Y\cap \maximalTightSetContainingHyperedge{p, H, w}(e) = \emptyset$ by \Cref{lem:weak-to-strong:uncrossing:maximal-tight-set-family-disjoint}, contradicting $\emptyset\neq e\cap X\subseteq e \cap Y \subseteq \maximalTightSetContainingHyperedge{p, H, w}(e)\cap Y$. 

    \item Follows from the definition of \merge operation (see \Cref{def:(pHw)-tight-sets}).
    
    \item 
    %We recall that the weight function $w^M = w - \alpha^M\chi_e - \alpha^M \chi_f + \alpha^M\chi_{e\cup f}$. 
    Let $X \subseteq V$. If at most one of the sets $X\cap e$ or $X\cap f$ is non-empty, then by part (2) of the current lemma (shown above) we have that $b_{(H^M, w^M)}(X) = b_{(H,w)}(X) \geq p(X)$. Suppose that $e\cap X, f\cap X \not = \emptyset$. Then, we have that $\alpha^M \leq \beta^M \leq b_{(H,w)}(X) - p(X)$.
    %\min\left\{b_{(H,w)}(X) - p(X) : X \subseteq V \text{ and } e \cap X, f \cap X \not = \emptyset\right\}$. 
    Thus,  
    $b_{(H^M, w^M)}(X) = b_{(H,w)}(X) - \alpha^M \geq b_{(H,w)}(X) - \beta^M \geq p(X)$, where the equality is by part (2) of the current lemma (shown above).

    \item We note that $w^M(e\cup f) = \alpha^M \in \Z_+$ by parts (1) and (2) of the current lemma (shown above). Furthermore, $b_{H^M,w^M}(X) \geq p(X)$ for every $X\subseteq V$ by part (3) of the current lemma (shown above). Thus, $\beta^R \in \Z_{\geq 0}$ since $p$ is an integer-valued function.

    \item Follows from the definition of \reduce operation (see \Cref{def:(pHw)-tight-sets}).

    \item Let $X\subseteq V$. First, suppose that the hyperedge $e\cup f$ is not contained in the set $X$. 
    %Let $(H_0, w_0)$ be the hypergraph defined by \Cref{alg:WeakToStrongCover}. 
    Then, we have that $b_{(H^R, w^R)}(X) = b_{(H^M, w^M)}(X) - d_{(H_0, w_0)}(X) \geq p(X) - d_{(H_0, w_0)}(X) = p^R(X)$. Here, the first equality is by part (5) of the current lemma and the inequality is by part (3) of the current lemma (proofs of both parts shown above).  Next, suppose that $e\cup f \subseteq X$. Then, we have that $b_{(H^R, w^R)}(X) = b_{(H^M, w^M)}(X) - \alpha^R \geq b_{H^M,w^M}(X) - \beta^R \geq b_{(H^M, w^M)}(X) - (b_{(H^M, w^M)}(X) - p(X)) = p(X) = p^R(X)$, where the first equality is by part (5) of the current lemma (shown above) and because $e\cup f \subseteq X$, and the final equality is because $e\cup f \subseteq X$.
\iffalse
    \item We prove the claim by showing the implication in both directions.
    First, suppose that $\alpha^M = \beta^M$. Let $X \subseteq V$ witness $\beta^M$, i.e. $\beta^M = b_{(H,w)}(X) - p(X)$ and $e\cap X, f\cap X \not = \emptyset$. Then, we have that
    $b_{H^M,w^M}(X) = b_{(H,w)}(X) - \alpha^M = p(X)$, where the first equality is by part (2) of the claim (shown above). Thus, $\alpha^R \leq \beta^R \leq b_{H^M,w^M}(X) - p(X) = 0$ \shubhang{the second inequality is a bug? Why does $e\cup f \subseteq X$? -- this can be shown using lemma 8 and lemma 4.}\shubhang{todo: move this statement to the place that actually uses it (I think claim 14.2 is the only place -- need to verify.) and prove it there correctly as a claim} and consequently $\alpha^R = 0$ by part (4) of the claim (shown above). Next, suppose that $\alpha^R = 0$. Then, since $w^M(e\cup f) \geq \alpha^M > 0$ by part (1), we have that $\alpha^R = \beta^R = 0$. Thus, there exists a set $X \subseteq V$ such that $e\cup f \subseteq X$ and $b_{H,w^M}(X) = p(X)$. 
    %We note that by the algorithm's choice of hyperedges $e, f$, we have that $\maximalTightSetContainingHyperedge{p, H, w}(e) \not = \maximalTightSetContainingHyperedge{p, H, w}(f)$ and consequently $\maximalTightSetContainingHyperedge{p, H, w}(e)\cap \maximalTightSetContainingHyperedge{p, H, w}(f) = \emptyset$ by \Cref{lem:weak-to-strong:uncrossing:maximal-tight-set-family-disjoint}. Thus, there does not exist $Y \in \tightSetFamily{p, H, w}$ such that $e\cap Y, f\cap Y \not= \emptyset$. In particular, $X \not \in \tightSetFamily{p, H, w}$. 
     Then, we have the following: $$p(X) = b_{(H^M, w^M)}(X) = b_{(H,w)}(X) - \alpha^M \geq b_{(H,w)}(X) - \beta^M \geq b_{(H,w)}(X) - (b_{(H,w)}(X) - p(X)) = p(X).$$
    Thus, all inequalities are equations, and we have that $\alpha^M = \beta^M$.

    \item 
    %The claim follows from the construction of the hyperedge set $E'$.
    We recall that $E' = E^R = \supp({w^R})$ by the definition of the \reduce operation. Furthermore, $w^R = w - w_0^M + w_1^M - w_0$, where $w_0^M:\{e, f\}\rightarrow\{\alpha^M\}$, $w_1^M:\{e\cup f\} \rightarrow\{\alpha^M\}$, and $w_0:\{e\cup f\}\rightarrow\{\alpha^R\}$ as defined by \Cref{alg:WeakToStrongCover}. Thus, if $w(e) = \alpha^M$, then $w^R(e) = 0$ and consequently $e \not \in E'$. Using similar reasoning with hyperedges $f$ and $e\cup f$ gives the claim.
    \fi
\end{enumerate}
\end{proof}

With the previous properties established, we now show that the function $p'$ and hypergraph $(H', w')$ constructed by \Cref{alg:WeakToStrongCover}  as input to the subsequent recursive call of the algorithm satisfy the two hypothesis conditions of \Cref{thm:WeakToStrongCover:main}  --  the function $p'$ is symmetric skew-supermodular, and the hypergraph $(H', w')$ weakly covers the function $p'$. This is a necessary step in order to give an inductive proof of \Cref{thm:WeakToStrongCover:main}. 

\begin{lemma}\label{lem:weak-to-strong-cover:next-recursive-call-hypothesis-satification:symm-skew-supermod-strong-cover}
%Consider a recursive call of \Cref{alg:WeakToStrongCover}. 
Suppose that the input to \Cref{alg:WeakToStrongCover} is a non-empty hypergraph $(H = (V, E), w)$ and a symmetric skew-supermodular function $p:2^V\rightarrow \Z$ %satisfying conditions \ref{cond:WeakToStrong:(a)} and \ref{cond:WeakToStrong:(b)} 
such that every hyperedge $\Tilde{e} \in E$ is contained in some $(p, H, w)$-tight set, $b_{(H,w)}(u) \not = 0$ for every $u \in V$, and $b_{(H,w)}(X) \geq p(X)$ for every $X \subseteq V$. 
Let $p'$ and $(H' = (V', E'), w')$ be as defined by \Cref{alg:WeakToStrongCover}. Then, we have the following:
\begin{enumerate}
    \item the function $p'$ is symmetric skew-supermodular, and
    \item $b_{(H', w')}(X) \geq p'(X)$ for every $X \subseteq V'$.
\end{enumerate}
\end{lemma}
\begin{proof} We prove each claim separately below. Let $(H_0, w_0), p^R, \zeros,$ and $(H^R, w^R)$ be as defined by \Cref{alg:WeakToStrongCover}.
\begin{enumerate}
    \item We recall that $p^R =  p - d_{(H_0, w_0)}$. We first observe that the function $p^R$ is symmetric skew-supermodular because the function $p$ is symmetric skew-supermodular and the cut function $d_{(H_0, w_0)}$ is symmetric submodular. Next, we observe that the function $p' = \functionContract{p^R}{\zeros}$. The claim then follows because the contraction operation preserves skew-supermodularity and symmetry.
    \item Let $X \subseteq V'$. Let $R \subseteq \zeros$ be such that $p'(X) = \functionContract{p^R}{\zeros}(X) = p^R(X \cup R)$. Then, we have the following: $b_{(H', w')}(X) = b_{(H^R, w^R)}(X \cup R) \geq p^R(X \cup R) = p'(X)$. Here, the first equality is because $w' = w^R$ and $b_{(H^R, w^R)}(u) = 0$ for every $u \in R$ by the definition of $\zeros$. The inequality is by \Cref{lem:WeakToStrong:properties-of-alpha^M-and-alpha^R}(6).  
\end{enumerate}
\end{proof}

In order to show that the function $p'$ and $(H',w')$ constructed by \Cref{alg:WeakToStrongCover} satisfy conditions \ref{cond:WeakToStrong:(a)} and \ref{cond:WeakToStrong:(b)}, we examine how the families of $(p, H, w)$-tight sets change during a recursive call of the algorithm. 
%We recall that  a set $X\subseteq V$ is said to be \emph{tight} w.r.t. a function $p:2^V\rightarrow\Z$ and a hypergraph $(H = (V, E), w)$ if $p(X) = b_{(H,w)}(X)$. We also recall that $\tightSetFamily{p,H,w}$ denotes the family of tight sets w.r.t. the function $p$ and hypergraph $(H,w)$. 
The next lemma shows that $(p, H, w)$-tight sets persist through the merge and reduce steps of a recursive call. Furthermore, the lemma also shows that the projection of a $(p, H, w)$-tight set onto the ground set of the subsequent recursive call is  $(p', H', w')$-tight.

\begin{lemma}\label{lem:WeakToStrongCover:TightSetFamiliesDuringAlgorithm}
%Consider a recursive call of \Cref{alg:WeakToStrongCover}. Suppose that the input to the recursive call is a non-empty hypergraph $(H = (V, E), w)$ and a symmetric skew-supermodular function $p:2^V\rightarrow \Z$ such that $b_{(H,w)}(X) \geq p(X)$ for every $X \subseteq V$. Furthermore, suppose that $\maximalTightSetContainingHyperedge{p, H, w}(\Tilde{e}) \not = \emptyset$ for every hyperedge $\Tilde{e} \in E$ and $b_{(H,w)}(u) > 0$ for every vertex $u \in V$.
%Consider a recursive call of \Cref{alg:WeakToStrongCover}. 
Suppose that the input to \Cref{alg:WeakToStrongCover} is a non-empty hypergraph $(H = (V, E), w)$ and a symmetric skew-supermodular function $p:2^V\rightarrow \Z$ %satisfying conditions \ref{cond:WeakToStrong:(a)} and \ref{cond:WeakToStrong:(b)} 
such that every hyperedge $\Tilde{e} \in E$ is contained in some $(p, H, w)$-tight set, $b_{(H,w)}(u) \not = 0$ for every $u \in V$, and $b_{(H,w)}(X) \geq p(X)$ for every $X \subseteq V$. 
Let $(H^M, w^M), (H^R, w^R), p'$ 
and $(H' = (V', E'), w')$ be as defined by \Cref{alg:WeakToStrongCover}.
Then, we have the following:
     \begin{enumerate}
        \item $\tightSetFamily{p, H, w} \subseteq \tightSetFamily{p, H^M, w^M} \subseteq \tightSetFamily{p^R, H^R, w^R}$, and 
        \item $\tightSetFamily{p^R, H^R, w^R}|_{V'} \subseteq \tightSetFamily{p', H', w'}$.
    \end{enumerate}
\end{lemma}
\begin{proof} We prove the claims separately below. Let $e, f, \alpha^M, \zeros$ be as defined by \Cref{alg:WeakToStrongCover}.
\begin{enumerate}
    \item We will first show that $\tightSetFamily{p, H, w} \subseteq \tightSetFamily{p, H^M, w^M}$. Let $X \in \tightSetFamily{p, H, w}$. 
    %Here, we observe that $e\cap X, f\cap X \not = \emptyset$ as otherwise 
    Suppose $e\cap X\neq \emptyset$ and $f\cap X\neq \emptyset$, then 
        $b_{(H^M, w^M)}(X) = b_{(H,w)}(X) - \alpha^M = p(X) - \alpha^M < p(X)$, which is a contradiction. Here, the first equality is because $e\cap X\neq \emptyset$ and $f\cap X\neq \emptyset$ and the final inequality is by \Cref{lem:WeakToStrong:properties-of-alpha^M-and-alpha^R}(1). 
        So, at most one of $e\cap X$ and $f\cap X$ is non-empty and consequently, $b_{(H,w)}(X)=b_{(H^M, w^M)}(X)$. 
        Thus, we have that $p(X) = b_{(H,w)}(X) = b_{(H^M, w^M)}(X)$, i.e. $X \in \tightSetFamily{p, H^M, w^M}$. 

        We will now show the second inclusion in the claim, i.e. $\tightSetFamily{p, H^M, w^M} \subseteq \tightSetFamily{p^R, H^R, w^R}$. Let $X \in \tightSetFamily{p, H^M, w^M} $. First, suppose that $\alpha^R = 0$. Then, we have that $b_{(H^R, w^R)}(X) = b_{(H^M, w^M)}(X) = p(X) = p^R(X)$, where the first and final equalities are because $\alpha^R = 0$. Thus, $X \in \tightSetFamily{p^R, H^R, w^R}$. Next, suppose that $\alpha^R > 0$. We note that if $(e\cup f) \cap X = \emptyset$, then $b_{(H^R,w^R)}(X) = b_{(H^M,w^M)}(X) = p(X) = p^R(X)$ and so $X \in \tightSetFamily{p^R, H^R, w^R}$. Consider the case where $(e\cup f) \cap X \not = \emptyset$. Here, we may assume that 
        $e\cup f\in \delta(X)$ since otherwise, 
        we obtain the following contradiction  $b_{(H^R, w^R)}(X) = b_{(H^M, w^M)}(X) - \alpha^R = p(X) - \alpha^R < p(X) = p^R(X)$, where the final equality is because $d_{(H_0, w_0)}(X) = 0$ since $e\cup f \not \in \delta_H(X)$. 
        Consequently, we have that $b_{(H^R, w^R)}(X) = b_{(H^M, w^M)}(X) - \alpha^R = p(X) - \alpha^R = p^R(X)$, and the claim holds. 
        
        \item Let $X \in \tightSetFamily{p^R, H^R, w^R}$.  Let $\zeros$ be as defined by \Cref{alg:WeakToStrongCover}. Then, we have the following:
        $$p'(X - \zeros) = \functionContract{p^R}{\zeros}(X - \zeros) \geq p^R(X) = b_{(H^R, w^R)}(X) = b_{(H^R, w^R)}(X - \zeros) = b_{(H', w')}(X - \zeros),$$
        where the penultimate equality is because $\zeros = \{u \in V : b_{(H^R, w^R)}(u) = 0\}$. However, by \Cref{lem:weak-to-strong-cover:next-recursive-call-hypothesis-satification:symm-skew-supermod-strong-cover}, we have that $b_{H',w'}(X - \zeros) \geq p'(X - \zeros)$. Thus, the inequality in the above sequence is an equation and we have that $X - \zeros \in \tightSetFamily{p', H', w'}$.
\end{enumerate}
\end{proof}
%With the previous relationships between the families $\tightSetFamily{p, H, w}, \tightSetFamily{p, H, w^M}, \tightSetFamily{p^R, H, w^R}$ and $\tightSetFamily{p', H', w'}$ established, 
We now show another useful property of \Cref{alg:WeakToStrongCover} which we will need while reasoning about the convergence of \Cref{alg:WeakToStrongCover} in the next section.
\begin{lemma}\label{lem:WeakToStrong:alpha^R=0-iff-alpha^M=beta^M}
Suppose that the input to \Cref{alg:WeakToStrongCover} is a non-empty hypergraph $(H = (V, E), w)$ and a symmetric skew-supermodular function $p:2^V\rightarrow \Z$
such that every hyperedge $\Tilde{e} \in E$ is contained in some $(p, H, w)$-tight set, $b_{(H,w)}(u) \not = 0$ for every $u \in V$, and $b_{(H,w)}(X) \geq p(X)$ for every $X \subseteq V$. 
 Let $\alpha^M, \alpha^R$ and $\beta^M$ be as defined by \Cref{alg:WeakToStrongCover}.  
 Then, $\alpha^R = 0$ if and only if $\alpha^M = \beta^M$.
 %(i.e., $\alpha^R > 0$ if and only if $\alpha^M < \beta^M$) 
\end{lemma}
\begin{proof}
    We prove the claim by showing the implication in both directions. Let $e, f, (H^M, w^M), \beta^R$ be as defined by \Cref{alg:WeakToStrongCover}. 
    %\knote{Add all other quantities from the algorithm that are used in the proof to the above sentence.}
    First, suppose that $\alpha^M = \beta^M$. Let $X \subseteq V$ be such that $\beta^M = b_{(H,w)}(X) - p(X)$ and $e\cap X, f\cap X \not = \emptyset$ (such a set exists by definition of $\beta^M$). Then, we have that
    $b_{(H^M,w^M)}(X) = b_{(H,w)}(X) - \alpha^M = p(X)$, where the first equality is by \Cref{lem:WeakToStrong:properties-of-alpha^M-and-alpha^R}(2). Thus, $X \in \tightSetFamily{p, H^M, w^M}$. Let $Y \in \maximalTightSetFamily{p, H^M, w^M}$ such that $X \subseteq Y$. \Cref{claim:WeakToStrong:new-tight-set-containing-e-and-f} below says that the set $Y$ contains the hyperedge $e \cup f$. With this claim, we complete the proof of the forward direction as follows: By \Cref{claim:WeakToStrong:new-tight-set-containing-e-and-f}, we have that $\alpha^R \leq \beta^R \leq b_{H^M,w^M}(Y) - p(Y) = 0$, where the final equality is because $Y \in \tightSetFamily{p, H^M, w^M}$. Consequently, $\alpha^R = 0$ because $\alpha^R \geq 0$ by \Cref{lem:WeakToStrong:properties-of-alpha^M-and-alpha^R}(4). 
\begin{claim}\label{claim:WeakToStrong:new-tight-set-containing-e-and-f}
       $e \cup f \subseteq Y$.
    \end{claim}
    \begin{proof}
    By way of contradiction, say $(e\cup f) - Y \not = \emptyset$. Thus, $e - Y \not = \emptyset$ or $f - Y \not = \emptyset$ (or both). We consider the case where $e - Y \not = \emptyset$ and remark that the proof of the case where $f - Y \not = \emptyset$ is along the same lines.
        We recall that the hyperedge $e$ is contained in some $(p, H, w)$-tight set, and thus the set $\maximalTightSetContainingHyperedge{p, H, w}(e)$ is well-defined and non-empty. Furthermore, by \Cref{lem:WeakToStrongCover:TightSetFamiliesDuringAlgorithm}(1), we have that the set $\maximalTightSetContainingHyperedge{p, H, w}(e) \in \tightSetFamily{p, H^M, w^M}$. Consequently, there exists a set $Z \in \maximalTightSetFamily{p, H^M, w^M}$ such that $e \subseteq \maximalTightSetContainingHyperedge{p, H, w}(e) \subseteq Z$. Since $e - Y \not = \emptyset$ but $e \subseteq Z$, we have that $Y \not = Z$.  We recall that $e \cap X \not = \emptyset$ by our choice of the set $X$, and hence, $e \cap Y \not = \emptyset$. Furthermore, by \Cref{lem:WeakToStrong:properties-of-alpha^M-and-alpha^R}(3), we have that $b_{(H^M, w^M)}(A) \geq p(A)$ for every $A \subseteq V$. Thus, by \Cref{lem:weak-to-strong:uncrossing:uncrossing-lemma} applied to the function $p$, hypergraph $(H^M, w^M)$, sets $Y$ and $Z$, and the hyperedge $e$, we have that the set $Y \cup Z \in \maximalTightSetFamily{p, H^M, w^M}$, a contradiction to the maximality of the sets $Y, Z \in \maximalTightSetFamily{p, H^M, w^M}$.
    \end{proof}
    
    Next, we prove the reverse direction of the claim. Suppose that $\alpha^R = 0$. 
    %By way of contradiction, say $\alpha^M < \beta^M$. 
    We note that $w^M(e\cup f) \geq \alpha^M > 0$ by Lemma \ref{lem:WeakToStrong:properties-of-alpha^M-and-alpha^R}(2) and \ref{lem:WeakToStrong:properties-of-alpha^M-and-alpha^R}(1). Thus, we have that $\alpha^R < w^M(e\cup f)$ and so $\alpha^R = \beta^R$. It follows that $\beta^R=0$. Consequently, there exists a set $X \subseteq V$ such that $e\cup f \subseteq X$ and $b_{(H^M,w^M)}(X) = p(X)$ (by definition of $\beta^R$). 
     Then, we have the following: $$p(X) = b_{(H^M, w^M)}(X) = b_{(H,w)}(X) - \alpha^M \geq b_{(H,w)}(X) - \beta^M \geq b_{(H,w)}(X) - (b_{(H,w)}(X) - p(X)) = p(X).$$
    Thus, all inequalities are equations, and we have that $\alpha^M = \beta^M$. Here, the second equality is by \Cref{lem:WeakToStrong:properties-of-alpha^M-and-alpha^R}(2) and $e\cup f \subseteq X$. The second inequality is also because $e\cup f \subseteq X$.
\end{proof}

 We now show that if the inputs to \Cref{alg:WeakToStrongCover} satisfy conditions \ref{cond:WeakToStrong:(a)} and \ref{cond:WeakToStrong:(b)}, then the inputs to every subsequent recursive call of \Cref{alg:WeakToStrongCover} also satisfy the two conditions.

%\knote{$e$ is overloaded}\shubhang{changed to $\Tilde{e}$. Also added a comment at the start of the section about this.}
\begin{lemma}\label{lem:weak-to-strong-cover:next-recursive-call-hypothesis-satification:hyperedge-in-tight-set-positive-coverage}
%Consider a recursive call of \Cref{alg:WeakToStrongCover}. Suppose that the input to the recursive call is a non-empty hypergraph $(H = (V, E), w)$ and a symmetric skew-supermodular function $p:2^V\rightarrow \Z$ such that $b_{(H,w)}(X) \geq p(X)$ for every $X \subseteq V$. Furthermore, suppose that $\maximalTightSetContainingHyperedge{p, H, w}(\Tilde{e}) \not = \emptyset$ for every hyperedge $\Tilde{e} \in E$ and $b_{(H,w)}(u) > 0$ for every vertex $u \in V$. 
%Consider a recursive call of \Cref{alg:WeakToStrongCover}. 
Suppose that the input to \Cref{alg:WeakToStrongCover} is a non-empty hypergraph $(H = (V, E), w)$ and a symmetric skew-supermodular function $p:2^V\rightarrow \Z$ %satisfying conditions \ref{cond:WeakToStrong:(a)} and \ref{cond:WeakToStrong:(b)} 
such that every hyperedge $\Tilde{e} \in E$ is contained in some $(p, H, w)$-tight set, $b_{(H,w)}(u) \not = 0$ for every $u \in V$, and $b_{(H,w)}(X) \geq p(X)$ for every $X \subseteq V$. 
Let $(H' = (V', E'), w')$ and $p'$ be as defined by \Cref{alg:WeakToStrongCover}. Then, we have the following:
\begin{enumerate}
    \item every hyperedge $\Tilde{e} \in E'$ is contained in some $(p', H', w')$-tight set, and
%$\maximalTightSetContainingHyperedge{p'}(\tilde{e}) \not = \emptyset$ for every hyperedge $\tilde{e} \in E'$, and
    \item  $b_{H',w'}(u) > 0$ for every vertex $u \in V'$.
\end{enumerate}
%In particular, we have that the hypergraph $(H',w')$ and the function $p'$ satisfy conditions \ref{cond:WeakToStrong:(a)} and \ref{cond:WeakToStrong:(b)}.
\end{lemma}
\begin{proof} We prove the claims separately below. Let $e, f, \zeros, \alpha^R, \beta^R, p', (H^M, w^M), (H^R, w^R)$ be as defined by \Cref{alg:WeakToStrongCover}.
%\knote{Let ... be as defined by \Cref{alg:WeakToStrongCover}.}
\begin{enumerate}
    \item We consider two cases. First, suppose that $\Tilde{e} \in E \cap E'$. Since the hyperedge $\Tilde{e}$ is contained in some $(p, H, w)$-tight set, we have that $\maximalTightSetContainingHyperedge{p, H, w}(\Tilde{e})$ is well-defined and non-empty. Furthermore, by \Cref{lem:WeakToStrongCover:TightSetFamiliesDuringAlgorithm}(2), we have that $\maximalTightSetContainingHyperedge{p, H, w}(\Tilde{e})|_{\zeros} \in \tightSetFamily{p', H', w'}$. Consequently, $\Tilde{e} \subseteq \maximalTightSetContainingHyperedge{p, H, w}(\Tilde{e})|_{\zeros} \in \tightSetFamily{p', H', w'}$. Next, suppose that $\Tilde{e} \in E' - E$. In particular, this implies that $E' - E \not= \emptyset$. 
    We recall that $e, f \in E$ are the hyperedges chosen by  \Cref{alg:WeakToStrongCover}  \Cref{alg:WeakToStrongCover:(4)}. 
    By Lemma \ref{lem:WeakToStrong:properties-of-alpha^M-and-alpha^R}(5) and \ref{lem:WeakToStrong:properties-of-alpha^M-and-alpha^R}(2), we have that $\Tilde{e} = e\cup f$ and $\alpha^R = \beta^R < w^M(e \cup f)$. Consequently, $\zeros = \emptyset$ and $V' = V$. Let $Z \subseteq V$ be a set such that 
    %witnessing $\alpha^R = \beta^R$, i.e. 
    $\alpha^R = b_{(H^M, w^M)}(Z) - p(Z)$ and $e\cup f \subseteq Z$ (such a set exists by definition of $\beta^R$).
    %\knote{Avoid using this phrasing: Let .. be a set witnessing .... It is non-standard. Instead, just say, let ... be a set such that ``conditions'' hold (such a set exists by definition of $\beta^R$}\shubhang{changed. leaving the comment here as a reminder to change in future proofs.}. 
    Then, we have the following: 
    $$b_{H',w'}(Z) = b_{(H^R, w^R)}(Z) = b_{(H^M, w^M)}(Z) - \alpha^R = p(Z) = p^R(Z) = p'(Z),$$
    and consequently the set $Z \in \tightSetFamily{p', H', w'}$.
    Here, the first equality is by definitions of the weight function $w^R$ and the set $\zeros$. The second and the fourth equalities are because $e \cup f \subseteq Z$. The third equality is because the set $Z$ is a witness set for $\alpha^R = \beta^R$. 
    The fifth equality is because $\zeros = \emptyset$.  %Thus, there exists a set $Y \in \maximalTightSetFamily{p',H',w'}$ such that $e \cup f \subseteq Z \subseteq Y$ and consequently $\maximalTightSetContainingHyperedge{p'}(e \cup f) \not = \emptyset$.

    \item We recall that $V' = V - \zeros = V - \{u\in V : b_{(H^R, w^R)}(u) = 0\}$. The claim then follows because $w' = w^R$. 
    %\knote{$w'=w^R$.}.
    
\end{enumerate}
\end{proof}

We now show the first main result of the section which says that \Cref{alg:WeakToStrongCover} terminates, and the hypergraph returned by the algorithm indeed strongly covers the input symmetric skew-supermodular function, i.e. satisfies property \ref{thm:WeakToStrongCover:main:(1)} of \Cref{thm:WeakToStrongCover:main}.

\begin{lemma}\label{lem:weak-to-strong:correctness:main:strong-cover}
    %Consider a recursive call of \Cref{alg:WeakToStrongCover}. 
    Suppose that the input to \Cref{alg:WeakToStrongCover} is a  hypergraph $(H = (V, E), w)$ and a symmetric skew-supermodular function $p:2^V\rightarrow \Z$ %satisfying conditions \ref{cond:WeakToStrong:(a)} and \ref{cond:WeakToStrong:(b)} 
such that every hyperedge $\Tilde{e} \in E$ is contained in some $(p, H, w)$-tight set, $b_{(H,w)}(u) \not = 0$ for every $u \in V$, and $b_{(H,w)}(X) \geq p(X)$ for every $X \subseteq V$.  Then, \Cref{alg:WeakToStrongCover} terminates in finite number of recursive calls. Furthermore, let $(H^*, w^*)$ be the hypergraph returned by the algorithm. Then, $d_{(H^*, w^*)}(X) \geq p(X)$ for every $X\subseteq V$.
    % Then, we have the following:
    % \begin{enumerate}
    %     \item $d_{(H^*, w^*)}(X) \geq p(X)$ for every $X\subseteq V$, and
    %     \item the hypergraph $(H^*, w^*)$ is formed by merging hyperedges of the hypergraph $(H,w)$.
    % \end{enumerate}
\end{lemma}
\begin{proof}
%\Cref{lem:weak-to-strong-cover:next-recursive-call-hypothesis-satification:symm-skew-supermod-strong-cover} and \Cref{lem:weak-to-strong-cover:next-recursive-call-hypothesis-satification:hyperedge-in-tight-set-positive-coverage} 
By Lemmas \ref{lem:weak-to-strong-cover:next-recursive-call-hypothesis-satification:symm-skew-supermod-strong-cover} and \ref{lem:weak-to-strong-cover:next-recursive-call-hypothesis-satification:hyperedge-in-tight-set-positive-coverage}, the input to every recursive call (except the last one) of the execution satisfies the hypothesis of \Cref{cor:weak-to-strong:uncrossing:exist-two-hyperedges-contained-in-disjoint-maximal-tight-sets}, and consequently, during every recursive call of the execution (except the last one), there exist two hyperedges such that \Cref{alg:WeakToStrongCover} \Cref{alg:WeakToStrongCover:(4)} is well-defined.
We first show that 
%the execution is finite, i.e. the 
algorithm terminates in a finite number of recursive calls. 
Consider the potential function $\phi(H=(V, E), w) \coloneqq  \sum_{e \in E}w(e)$. We note that $\phi(H, w) \in \Z_{\geq 0}$ since the weight function $w:E\rightarrow\Z_+$ is a positive integral function. Furthermore, if $\phi(H, w) = 0$, then we have that $E = \emptyset$ and the algorithm terminates. We recall that $p'$ and $(H', w')$ denote the inputs to the subsequent recursive call of \Cref{alg:WeakToStrongCover}. 
By Lemmas \ref{lem:weak-to-strong-cover:next-recursive-call-hypothesis-satification:symm-skew-supermod-strong-cover} and \ref{lem:weak-to-strong-cover:next-recursive-call-hypothesis-satification:hyperedge-in-tight-set-positive-coverage}, we have that $p'$ is a symmetric skew-supermodular function and $(H' = (V', E'), w')$ is a non-empty hypergraph such that every hyperedge $\Tilde{e} \in E'$ is contained in some $(p', H', w')$-tight set, $b_{(H', w')}(u) \not = 0$ for every $u \in V'$, and $b_{(H', w')}(X) \geq p'(X)$ for every $X \subseteq V'$.
Then, by Lemma \ref{lem:WeakToStrong:properties-of-alpha^M-and-alpha^R}(1), \ref{lem:WeakToStrong:properties-of-alpha^M-and-alpha^R}(2) and \ref{lem:WeakToStrong:properties-of-alpha^M-and-alpha^R}(5), we have that $\phi(H, w) > \phi(H', w')$, i.e. the value of the potential function strictly decreases between subsequent recursive calls. Consequently, the algorithm terminates in a finite number of recursive calls.


We now show that the hypergraph $(H^*,w^*)$ returned by \Cref{alg:WeakToStrongCover} strongly covers the function $p$, i.e. $d_{(H^*, w^*)}(X) \geq p(X)$ for every $X\subseteq V$. We show this by induction on $w(E)$. For the base case, consider $w(E) = 0$. Since the weight function $w$ is positive, we have that $E = \emptyset$.
Let $X \subseteq V$. We have that $d_{(H, w)}(X) = 0 = b_{(H,w)}(X) \geq p(X)$ and so the claim holds. For the inductive case, consider $w(E) > 0$. Let $X\subseteq V$ be an arbitrary subset. 
Let $(H_0^* = (V', E_0^*), w_0^*)$ be the hypergraph returned by the subsequent recursive call on input hypergraph $(H' = (V', E'), w')$ and function $p'$. 
%Let $(\Tilde{H} = (V', \Tilde{E}), \Tilde{w})$ be the hypergraph returned by the subsequent recursive call on input hypergraph $(H' = (V', E'), w')$ and function $p'$. 
We recall that $w'(E') < w(E)$ and so by the inductive hypothesis, we have that $d_{(H_0^*, w_0^*)}(X) \geq p'(X)$ for every $X\subseteq V'$. We also recall that the hypergraph $(H^*, w^*) = (G+H_0, c + w_0)$, where the hypergraph $(G, c)$ is obtained from $(H_0^*, w_0^*)$ by adding the vertices $\zeros$. Then, we have the following:
\begin{align*}
    d_{(H^*, w^*)}(X)& = d_{G,c}(X) + d_{(H_0, w_0)}(X)&\\
    & = d_{H_0^*, w_0^*}(X) + \alpha^R\indicator_{e\cup f\in \delta(X)}&\\
    &\geq p'(X) + \alpha^R\indicator_{e\cup f\in \delta(X)}& \text{(by inductive hypothesis)}\\
    & = \functionContract{\left(p - \alpha^R\cdot \indicator_{e\cup f\in \delta(X)}\right)}{\zeros}(X) + \alpha^R\indicator_{e\cup f\in \delta(X)} &\\
    &\geq \left(p - \alpha^R\cdot  \indicator_{e\cup f\in \delta(X)}\right)(X) + \alpha^R\indicator_{e\cup f\in \delta(X)}&\\
    &= p(X).&
\end{align*}
\end{proof}

We now show that the hypergraph returned by \Cref{alg:WeakToStrongCover} is obtained by merging hyperedges of the input hypergraph (see \Cref{def:WeakToStrongCover:hyperedge-merge} for the definition of merging hyperedges). 
This will imply that the hypergraph returned by the algorithm satisfies property \ref{thm:WeakToStrongCover:main:(2)} of \Cref{thm:WeakToStrongCover:main}. To show this property, we will need the following lemma about the choice of hyperedges $e$ and $f$ being merged by the algorithm.

\begin{lemma}\label{lem:WeakToStrongCover:properties-of-e-and-f-chosen-during-alg}
    %Consider a recursive call of \Cref{alg:WeakToStrongCover}. 
    Suppose that the input \Cref{alg:WeakToStrongCover} call is a non-empty hypergraph $(H = (V, E), w)$ and a symmetric skew-supermodular function $p:2^V\rightarrow \Z$ %satisfying conditions \ref{cond:WeakToStrong:(a)} and \ref{cond:WeakToStrong:(b)} 
such that every hyperedge $\Tilde{e} \in E$ is contained in some $(p, H, w)$-tight set, $b_{(H,w)}(u) \not = 0$ for every $u \in V$, and $b_{(H,w)}(X) \geq p(X)$ for every $X \subseteq V$.  Let $e, f$ be as defined by \Cref{alg:WeakToStrongCover}. Then,
\begin{enumerate}
    \item $e \cap f = \emptyset$, and
    \item $e\cup f \not \in E$.
\end{enumerate}
\end{lemma}
\begin{proof} We prove the claims separately below.
    \begin{enumerate}
        \item By the choice of hyperedges $e$ and $f$, we have that $\maximalTightSetContainingHyperedge{p, H, w}(e) \not = \maximalTightSetContainingHyperedge{p, H, w}(f)$. We observe that by maximality of the sets $\maximalTightSetContainingHyperedge{p, H, w}(e)$ and $\maximalTightSetContainingHyperedge{p, H, w}(f)$, we have that $\maximalTightSetContainingHyperedge{p, H, w}(e), \maximalTightSetContainingHyperedge{p, H, w}(f) \in \maximalTightSetFamily{p, H, w}$. By \Cref{lem:weak-to-strong:uncrossing:maximal-tight-set-family-disjoint}, we have that $\maximalTightSetContainingHyperedge{p, H, w}(e) \cap  \maximalTightSetContainingHyperedge{p, H, w}(f) = \emptyset$. Consequently, $e \cap f = \emptyset$.
        \item By way of contradiction, suppose that $e\cup f \in E$. Thus, $\maximalTightSetContainingHyperedge{p, H, w}(e\cup f) \not = \emptyset$. Since $e, f \subseteq \maximalTightSetContainingHyperedge{p, H, w}(e\cup f)$, we have that $\maximalTightSetContainingHyperedge{p, H, w}(e) = \maximalTightSetContainingHyperedge{p, H, w}(f) = \maximalTightSetContainingHyperedge{p, H, w}(e\cup f)$ by the maximality of the three sets, contradicting the choice of hyperedges $e, f$ made by the algorithm.
    \end{enumerate}
\end{proof}

We now conclude the section by showing that the hypergraph returned by \Cref{alg:WeakToStrongCover} indeed satisfies property \ref{thm:WeakToStrongCover:main:(2)} of \Cref{thm:WeakToStrongCover:main}.

\begin{lemma}\label{lem:weak-to-strong:correctness:main:hyperedge-merge}
    %Consider a recursive call of \Cref{alg:WeakToStrongCover}. 
    Suppose that the input to \Cref{alg:WeakToStrongCover} is a hypergraph $(H = (V, E), w)$ and a symmetric skew-supermodular function $p:2^V\rightarrow \Z$ %satisfying conditions \ref{cond:WeakToStrong:(a)} and \ref{cond:WeakToStrong:(b)} 
such that every hyperedge $\Tilde{e} \in E$ is contained in some $(p, H, w)$-tight set, $b_{(H,w)}(u) \not = 0$ for every $u \in V$, and $b_{(H,w)}(X) \geq p(X)$ for every $X \subseteq V$. 
    Then, the hypergraph $(H^* = (V, E^*), w^*)$ returned by the algorithm is obtained by merging hyperedges of the hypergraph $(H, w)$.
\end{lemma}
\begin{proof}
     We show the claim by induction on $w(E)$. For the base case, consider $w(E) = 0$. Since the weight function $w$ is positive, we have that $E = \emptyset$ and hence $(H, w)$ is the empty hypergraph. In this case, the hypergraph $(H^*, w^*)$ returned by the algorithm is also the empty hypergraph and so the claim holds. 

    For the inductive case, suppose that $w(E) > 0$. Let $(H^M, w^M), \alpha^M$, $(H' = (V', E'), w'), p',(H^R, w^R)$, $(H_0 = (V, E_0), w_0), \zeros, (H^*_0, w^*_0), (G, c)$ be as defined by \Cref{alg:WeakToStrongCover}. By Lemmas \ref{lem:weak-to-strong-cover:next-recursive-call-hypothesis-satification:symm-skew-supermod-strong-cover} and \ref{lem:weak-to-strong-cover:next-recursive-call-hypothesis-satification:hyperedge-in-tight-set-positive-coverage}, the inputs $(H',w')$ and $p'$ satisfy the hypothesis of \Cref{lem:weak-to-strong:correctness:main:strong-cover}. 
    Thus, the execution of the algorithm from this recursive call terminates in a finite number of recursive calls and returns the hypergraph $(H_0^* = (V', E_0^*), w_0^*)$. 
    %The following two claims show that the hypergraphs maintained and returned by the algorithm can be obtained by merging hyperedges of appropriate hypergraphs.
    Then, by Claims \ref{claim:WeakToStrong:MergeProperty:H^*w^*-from-H^Mw^M} and \ref{claim:WeakToStrong:MergeProperty:H^Mw^M-from-Hw} below, we have that the hypergraph $(H^*, w^*)$ is obtained by merging hyperedges of $(H, w)$.
    \end{proof}
    \iffalse
\begin{claim}\label{claim:WeakToStrong:Merge-Property:Gc-from-H^Rw^R}
        $(G, c)$ is formed by merging hyperedges of $(H^R, w^R)$.
    \end{claim}
    \begin{proof}
        By Lemma \ref{lem:WeakToStrong:properties-of-alpha^M-and-alpha^R}(1), \ref{lem:WeakToStrong:properties-of-alpha^M-and-alpha^R}(2) and \ref{lem:WeakToStrong:properties-of-alpha^M-and-alpha^R}(5), we have that $w'(E') < w(E)$. Thus, by the inductive hypothesis, we have that the hypergraph $(H_0^*, w_0^*)$ is formed by merging hyperedges of $(H', w')$. The claim then follows because of the following two observations: $(H', w')$ is obtained from $(H^R, w^R)$ by removing the vertices $\zeros$ and $(G, c)$ is formed by adding vertices $\zeros$ to $(H_0^*, w_0^*)$.
    \end{proof}
\fi
    \begin{claim}\label{claim:WeakToStrong:MergeProperty:H^*w^*-from-H^Mw^M}
        $(H^*, w^*)$ is obtained by merging hyperedges of $(H^M, w^M)$.
    \end{claim}
    \begin{proof}
        We first show that the hypergraph $(G, c)$ is obtained by merging hyperedges of $(H^R, w^R)$.
         By Lemma \ref{lem:WeakToStrong:properties-of-alpha^M-and-alpha^R}(1), (2), and (5), we have that $w'(E') < w(E)$. Thus, by the inductive hypothesis, we have that the hypergraph $(H_0^*, w_0^*)$ is obtained by merging hyperedges of $(H', w')$. Consequently, $(G, c)$ is formed by merging hyperedges of $(H^R, w^R)$ because of the following two observations: $(H', w')$ is obtained from $(H^R, w^R)$ by removing the vertices $\zeros$ and $(G, c)$ is obtained by adding vertices $\zeros$ to $(H_0^*, w_0^*)$.
        
        Next, we recall that $(H^*, w^*) = (G + H_0, c + w_0)$. Thus, by the conclusion of the previous paragraph, we have that $(H^*, w^*)$ is obtained by merging hyperedges of $(H^R + H_0, w^R + w_0)$. The claim then follows because $(H^R + H_0, w^R + w_0) = (H^M, w^M)$ by \Cref{lem:WeakToStrong:properties-of-alpha^M-and-alpha^R}(5).
    \end{proof}

    \begin{claim}\label{claim:WeakToStrong:MergeProperty:H^Mw^M-from-Hw}
        $(H^M,w^M)$ is obtained by merging hyperedges of $(H, w)$.
    \end{claim}
    \begin{proof}
        We recall that $(H^M, w^M) = \merge((H,w), e, f, \alpha^M)$ (see \Cref{alg:WeakToStrongCover} \Cref{alg:WeakToStrongCover:(6)}).
        Thus, by definition of the \merge operation, the hypergraph $(H^M, w^M)$ is obtained by decreasing the weight of hyperedges $e$ and $f$ by $\alpha^M$ and increasing the weight of the hyperedge $e\cup f$ by $\alpha^M$. Consequently, the multi-hypergraph $H^M_{w^M}$ can be obtained from the multi-hypergraph $H_w$ by repeatedly merging hyperedges $e$ and $f$ exactly $\alpha^M$ times. We note that the hyperedges $e$ and $f$ are disjoint by \Cref{lem:WeakToStrongCover:properties-of-e-and-f-chosen-during-alg}(1). Thus, the hypergraph $(H^M,w^M)$ is obtained by merging hyperedges of $(H, w)$ by definition. 
        %We recall that $H_0 = (V, E_0 = \{e\cup f\})$. Let $m_0 : E_0 \rightarrow \{\alpha^M\}$. Furthermore,  let $H_1 := (V, E_1:=\{e, f\})$ and $m_1: E_1 \rightarrow \{\alpha^M\}$. By \Cref{lem:WeakToStrong:properties-of-alpha^M-and-alpha^R}(2), we have that $(H^M, w^M) = (H + H_0 - H_1, w + m_0 - m_1)$ \shubhang{Is the ``hypergraph subtraction'' notation ok to use?} \knote{No, don't use subtraction notation. It raises confusion about whether the hyperedge weights are negative.}. We note that the hypergraph $(H_0, m_0)$ can be formed by merging hyperedges of $(H_1, m_1)$  since the hyperedges $e$ and $f$ are disjoint by \Cref{lem:WeakToStrongCover:properties-of-e-and-f-chosen-during-alg}(1).
    \end{proof}

    %By Claims \ref{claim:WeakToStrong:MergeProperty:H^*w^*-from-H^Mw^M} and \ref{claim:WeakToStrong:MergeProperty:H^*w^*-from-H^Mw^M} above, we have that the hypergraph $(H^*, w^*)$ is formed from the hypergraph $(H, w)$ by merging hyperedges.

    \iffalse
    For notational convenience, we use  $(G^* = (V', F^*), c^*)$ to denote the  hypergraph $ (H_0^* = (V', E_0^*), w_0^*)$. By Lemma \ref{lem:WeakToStrong:properties-of-alpha^M-and-alpha^R}(1), \ref{lem:WeakToStrong:properties-of-alpha^M-and-alpha^R}(2) and \ref{lem:WeakToStrong:properties-of-alpha^M-and-alpha^R}(5), we have that $w'(E') < w(E)$. Thus, by the inductive hypothesis, we have that the hypergraph $(G^*, c^*)$ is formed by merging hyperedges of the hypergraph $(H', w')$. In particular, the multi-hypergraph $G^*_{c^*}$ is obtained from the multi-hypergraph $H'_{w'}$ by repeatedly merging two disjoint hyperedges. Consider a sequence of $k \in \Z_{\geq 0}$ operations of merging disjoint hyperedges starting from the multi-hypergraph $H'_{w'}$ and resulting in the multi-hypergraph $G^*_{c^*}$. Let $\sigma = (e_1, f_1), (e_2, f_2), \ldots, (e_k, f_k)$ denote the collection of pairs of (disjoint) hyperedges merged in order in this sequence. We now use $\sigma$ to construct a new collection $\sigma'$ such that starting from the multi-hypergraph $H_w$ and sequentially merging the pairs of hyperedges in $\sigma'$   results in the multi-hypergraph $H^*_{w^*}$. We define the collection $\sigma'$ in two steps as follows.
    \begin{enumerate}
        \item Consider the sequence $\sigma$. For every $i \in [k]$ such that $e\cup f \in \{e_i, f_i\}$, add the pair $(e, f)$ before the pair $(e_i, f_i)$. Furthermore, let $q := |\{i \in [k] : e\cup f \in \{e_i, f_i\}\}|$; append $q$ pairs of $(e, f)$ at the end of the previous sequence. Denote the resulting sequence as $\sigma_1'$.
        \item Let $p := \alpha^M - \alpha^R$. Then, we define the sequence $\sigma_1'$ with an additional $p$ pairs of $(e, f)$ appended at the end.
    \end{enumerate}
    
    We recall that by Lemma \ref{lem:WeakToStrong:properties-of-alpha^M-and-alpha^R}(2) and \ref{lem:WeakToStrong:properties-of-alpha^M-and-alpha^R}(5), we have that $w'(e\cup f) = \alpha^M - \alpha^R = p$. Consequently, there are $p$ copies of the hyperedge $e\cup f$ in the multi-hypergraph $H'_{w'}$. 
\shubhang{todo: new proof below}
    \begin{claim}
        The multi-hypergraph $G_c$ can be obtained from the multi-hypergraph $H'_{w'}$ by merging hyperedges.
    \end{claim}
    \begin{proof}
        By the inductive hypothesis, the hypergraph $(H^*_0, w^*_0)$ can be obtained from the hypergraph $(H', w')$ by merging hyperedges. Since the hypergraph $(G, c)$ is formed from $(H_0^*, w_0^*)$ by adding vertices $\zeros$, it follows that the hypergraph $(G, c)$ can be obtained from the hypergraph $(H', w')$ by merging hyperedges.            
    \end{proof}

    Let $(\Tilde{H}_0 := (V, \Tilde{E}_0 :=\{e, f\})$ and $\Tilde{w}_0 : \Tilde{E_0} \rightarrow\{ \alpha^R\}\}$ 
    \begin{claim}
        The multi-hypergraph $H_{0_{w_0}}$ can be obtained from the multi-hypergraph $\Tilde{H}_{0_{\Tilde{w}_0}}$.
    \end{claim}
    \begin{proof}
        \shubhang{todo}
    \end{proof}

    \begin{claim}
        $H_w = G_c + \Tilde{H}_{0_{\Tilde{w}_0}}$.
    \end{claim}
    \begin{proof}
        \shubhang{todo}.
    \end{proof}

\shubhang{end new proof}
    
\shubhang{old proof below}

    
    $(G^* = (V', F^*), c^*)$. Furthermore, the inputs 
$(H', w')$ and $p'$ also satisfy the hypothesis of the claim, and $w'(E') < w(E)$ by \Cref{lem:WeakToStrong:properties-of-alpha^M-and-alpha^R}(2), \Cref{lem:WeakToStrong:properties-of-alpha^M-and-alpha^R}(5) and \Cref{lem:WeakToStrong:properties-of-alpha^M-and-alpha^R}(7). Thus, by the inductive hypothesis, we have that the hypergraph $(G^*, c^*)$ is formed by merging hyperedges of the hypergraph $(H', w')$. In particular, the multi-hypergraph $G^*_{c^*}$ is obtained from the multi-hypergraph $H'_{w'}$ by repeatedly merging two disjoint hyperedges. Consider a sequence of operations of merging disjoint hyperedges starting from the hypergraph $H'_{w'}$ and resulting in the hypergraph $G^*_{c^*}$. We now use this sequence to construct a new sequence of merge operations such that starting from the multi-hypergraph $H_w$ and performing the operations along the sequence results in the multi-hypergraph $H^*_{w^*}$. For this, we start with the multihypergraph $H_w$ and consider the aforementioned sequence guaranteed by the inductive hypothesis  --  we note that this may not a valid sequence of merge operations for the mutihypergraph $H_w$ as the hyperedge $e\cup f \not \in E$ by \Cref{lem:WeakToStrongCover:properties-of-e-and-f-chosen-during-alg}(2), but the hyperedge $e\cup f$ may be used in a merge operation in the sequence. By \Cref{lem:WeakToStrong:properties-of-alpha^M-and-alpha^R}(2) and \Cref{lem:WeakToStrong:properties-of-alpha^M-and-alpha^R}(5), we have that $w' = w^R = w - w_0^M + w_1^M - w_0$, where $w_0^M:\{e, f\}\rightarrow\{\alpha^M\}$, $w_1^M:\{e\cup f\} \rightarrow\{\alpha^M\}$, and $w_0:\{e\cup f\}\rightarrow\{\alpha^R\}$ as defined by \Cref{alg:WeakToStrongCover}. Consequently there are $\alpha^M - \alpha^R$ copies of the hyperedge $e\cup f$ in the multihypergraph $H'_{w'}$. Suppose that $k$ of these copies are used during merge operations along the sequence. We preface each of these $k$ operations with the operation of merging hyperedges $e$ and $f$. Furthermore, at the end of the sequence, we include $\alpha^M - \alpha^R - k$ new operations of merging hyperedges $e$ and $f$. By \Cref{lem:WeakToStrongCover:properties-of-e-and-f-chosen-during-alg}(1), the hyperedges $e$ and $f$ are disjoint. Thus, the resulting sequence of merge operations becomes a valid sequence of merge operations starting from the multihypergraph $H_w$. Furthermore, this new sequence of merge operations results in the multihypergraph $G^*_{c^*}$ along with an additional $\alpha^R$ copies of the hyperedges $e$ and $f$. By \Cref{alg:WeakToStrongCover} \Cref{alg:WeakToStrongCover:(11)}, the hypergraph $(H^*, w^*) = (G^* + H_0, c^* + w_0)$. Thus, we append the sequence with $\alpha^R$ new operations of merging the hyperedges $e$ and $f$ to get the multihypergraph $H^*_{w^*}$ as desired.
\fi
%\end{proof}

%% proof using induction on total weight and old definition of merging
\iffalse
\begin{proof}
    We show the claim by induction on $w(E)$. For the base case, consider $w(E) = 0$. Since the weight function $w$ is positive, we have that $E = \emptyset$ and hence $(H, w)$ is the empty hypergraph. In this case, the hypergraph $(H^*, w^*)$ returned by the algorithm is also the empty hypergraph and so the claim holds. 

    For the inductive case, suppose that $w(E) > 0$. Let the hypergraph $(H' = (V', E'), w')$ and the function $p':2^{V'}\rightarrow\Z$ be the inputs to the subsequent recursive call of the algorithm. By \Cref{lem:weak-to-strong-cover:next-recursive-call-hypothesis-satification:symm-skew-supermod-strong-cover} and \Cref{lem:weak-to-strong-cover:next-recursive-call-hypothesis-satification:hyperedge-in-tight-set-positive-coverage}, the inputs $(H',w')$ and $p'$ satisfy the hypothesis of \Cref{lem:weak-to-strong:correctness:main:strong-cover}. Thus, the the execution of the algorithm from this recursive call terminates in finite time and returns a hypergraph $(G^* = (V', F^*), c^*:F^*\rightarrow\Z_+)$. Furthermore, the inputs 
$(H', w')$ and $p'$ also satisfy the hypothesis of the claim, and $w'(E') < w(E)$ by \Cref{lem:WeakToStrong:properties-of-alpha^M-and-alpha^R} part (6). By the inductive hypothesis, we have that the hypergraph $(G^*, c^*)$ is formed by merging hyperedges of the hypergraph $(H', w')$. We now make some assumptions that will simplify the exposition of the proof. We henceforth assume that  $e, f \in E'$. We can make this assumption because if the hyperedge $e \not \in E'$, then we can add the hyperedge into $E'$ and set $w'(e) \coloneqq  0$. We similarly handle the case when the hyperedge $f \not \in E'$. We also assume that the hyperedge $e \cup f \in F^*$ and $e \cup f \in E$. As in the explanation of the previous assumption, if $e\cup f \not \in F^*$, then we add the hyperedge $e \cup f$ to $F^*$ and set $c^*(e\cup f) \coloneqq  0$. We note that $e\cup f \not \in E$ and so we add the hyperedge to $E$ with $w(e\cup f) \coloneqq  0$. With these assumption and by \Cref{lem:WeakToStrong:properties-of-alpha^M-and-alpha^R}, we have that $E = E'$. Furthermore, by the algorithm description, we have that  $E^* = F^*$. \knote{These assumptions are confusing to track inside the claims. For example, claim 42.1 is already confusing. Can the proof be written without making these assumptions? Perhaps, use indicator functions if you need to condition? Or, it is okay if you want to extend the domain of $\tilde{c}$ to $E^*\times E$ by defining its values to be zero on the new elements of the domain. }

    %Let $e^* \in E^*$ be an arbitrary hyperedge.
    By the inductive hypothesis, there exists a function $\Tilde{c}:F^*\times E' \rightarrow \Z_{\geq 0}$ satisfying property (1) of \Cref{def:WeakToStrongCover:hyperedge-merge}, and for every hyperedge $e^*\in E^*$, there exist positive integers $k' \in \Z_+$ and $\gamma_1', \ldots, \gamma_{k'}' \in \Z_+$, and a collection of hyperedge subsets $S_1', \ldots S_{k'}' \subseteq E'$  satisfying properties (2a)-(2c) of \Cref{def:WeakToStrongCover:hyperedge-merge}.

    We now extend these  quantities to simplify the exposition of the proof. First, suppose that the hyperedge $e \not \in E'$. Then, we extend the function $\Tilde{c}$ by defining $\Tilde{c}(e^*, e) \coloneqq  0$ for every $e^* \in F^* \cup \{e\cup f\}$. We perform a similar extension in the case where the hyperedge $f \not \in E'$. Next, suppose that the hyperedge $e \cup f \not \in F^*$. We extend the function $\Tilde{c}$ by defining the function value $\Tilde{c}(e\cup f, \Tilde{e}) \coloneqq  0$ for every $\Tilde{e} \in E' \cup \{e, f\}$. Here, we also define $k' \coloneqq  0$. We emphasize that these extensions are merely for notational convenience.
   
    We now define the weight function $\Tilde{w}:E^*\times E \rightarrow \Z_{\geq 0}$ as follows. Let $\alpha^R$ be the quantity defined by the recursive call of \Cref{alg:WeakToStrongCover}.
    %Let $e^* \in E_{i+1}^*$, and let $k', \gamma_1', \ldots, \gamma_{k'}' \in \Z_+$, and  $S_1', \ldots S_{k'}' \subseteq E_{i+1}$ be  values for $e^*$ satisfying (\ref{tag:hyperedge-merge-properties}). 
    For every $e^*\in E^*$ and $\Tilde{e} \in E$, we let



$$\Tilde{w}(e^*, \Tilde{e}) \coloneqq  \begin{cases}
        \Tilde{c}(e^*, \Tilde{e})&\text{ if } \Tilde{e} \not\in \{e, f\}\\
        \Tilde{c}(e^*, \Tilde{e}) + \Tilde{c}(e^*, e \cup f)&\text{ if } \Tilde{e} \in \{e, f\} \text{ and } e^* \not = e\cup f\\
        \Tilde{c}(e^*, \Tilde{e}) + \alpha^R& \text{ otherwise (i.e., if $\Tilde{e} \in \{e, f\}$ and $e^* = e \cup f$).}
    \end{cases}$$

    \Cref{claim:weak-to-strong:correctness:main:hyperedge-merge:property1} below says that the weight function $\Tilde{w}$ satisfies property (1) of \Cref{def:WeakToStrongCover:hyperedge-merge}. For the hyperedge $e^*$, we now define positive integers $k \in \Z_+$ and $\gamma_1, \ldots, \gamma_{k} \in \Z_+$, and the collections of hyperedge subsets $S_1, \ldots S_{k} \subseteq E$ as follows:
    \begin{align*}
        k &: =\begin{cases}
            k' + 1& \text{ if } e^* = e\cup f\\
            k'&\text{ otherwise.}
        \end{cases}&\\
        \gamma_j&\coloneqq  \begin{cases}
            \alpha^R& \text{ if } e^* = e\cup f \text{ and } j = k\\
            \gamma_j'&\text{ otherwise}.
        \end{cases}&\\
        S_j& \coloneqq  \begin{cases}
            \{e, f\}&\text{ if } e^* = e\cup f \text{ and } j = k\\
            S_j' - \{e\cup f\} + \{e, f\}&\text{ if } e\cup f \in S_j'\\
            S_j'& \text{ otherwise.}
        \end{cases}&
    \end{align*}
    \knote{Should the conditions for $\gamma_j$ and $S_j$ be for $j=k+1$?} 
    \Cref{claim:weak-to-strong:correctness:main:hyperedge-merge:property2}, \Cref{claim:weak-to-strong:correctness:main:hyperedge-merge:property3} and \Cref{claim:weak-to-strong:correctness:main:hyperedge-merge:property4} below show that the function $\Tilde{w}$, and the quantities $k$, $\gamma_1, \ldots, \gamma_k$ and $S_1, \ldots, S_k$ satisfy properties (2a)-(2c) of \Cref{def:WeakToStrongCover:hyperedge-merge} for the hyperedge $e^*$ and thus complete the proof of the claim. \knote{Phrase the claims to resemble Definition 4. }
\end{proof}

    \begin{claim}\label{claim:weak-to-strong:correctness:main:hyperedge-merge:property1}
        For every $\Tilde{e} \in E$, we have that $\sum_{e^* \in E^*}\Tilde{w}(e^*, \Tilde{e}) = w(\Tilde{e})$.
    \end{claim}
    \begin{proof}
        We note that $e \cup f \not \in E$ \knote{$w(e\cup f)=0$?} by \Cref{lem:WeakToStrongCover:properties-of-e-and-f-chosen-during-alg}. Let $\Tilde{e} \in E$. First, suppose that $\Tilde{e} \not \in \{e, f\}$. Then, we have the following:
        $$\sum_{e^* \in E^*}\Tilde{w}(e^*, \Tilde{e}) = \sum_{e^* \in F^*}\Tilde{c}(e^*, \Tilde{e}) = w'(\Tilde{e}) = w(\Tilde{e}).$$
        Here, the first and final equalities are because $\Tilde{e} \not \in \{e, f, e\cup f\}$ and the second equality is by the inductive hypothesis. Next, suppose that $\Tilde{e} \in \{e, f\}$. Then, we have the following:
        \begin{align*}
            \sum_{e^* \in E^*}\Tilde{w}(e^*, \Tilde{e})&= \sum_{e^* \in F^*}\Tilde{c}(e^*, \Tilde{e}) + \sum_{e^* \in F^*}\Tilde{c}(e^*, e\cup f) + \alpha^R&\\
            &= w'(\Tilde{e}) + w'(e\cup f) + \alpha^R&\\
            & = w(\Tilde{e}) - \alpha^M + (\alpha^M - \alpha^R) + \alpha^R&\\
            &= w(\Tilde{e}).  &
        \end{align*}
        Here, the second equality is by the inductive hypothesis and the third equality is by \Cref{lem:WeakToStrong:properties-of-alpha^M-and-alpha^R}.
    \end{proof}

\begin{claim}\label{claim:weak-to-strong:correctness:main:hyperedge-merge:property2}
    For every $j \in [k]$, we have that $e^* = \biguplus_{\Tilde{e} \in S_j}\Tilde{e}$.
\end{claim}
\begin{proof}
    We consider three cases based on the definition of the collection $S_j$. First, suppose that $e^* = e \cup f$ and $j = k$. Then, we have that $S_j = \{e, f\}$. We recall that the hyperedges $e$ and $f$ are disjoint by \Cref{lem:WeakToStrongCover:properties-of-e-and-f-chosen-during-alg} and so the claim holds. Next, suppose that $e \cup f \in S_j'$. Then, we have that $\uplus_{\Tilde{e} \in S_j} \Tilde{e} = \uplus_{\Tilde{e} \in S_j'} \Tilde{e} = e^*$, where the second equality is  because the hyperedges $e$ and $f$ are disjoint by \Cref{lem:WeakToStrongCover:properties-of-e-and-f-chosen-during-alg}, and the final equality is by the inductive hypothesis. Finally, suppose that $e\cup f \not \in S_j'$. Then, we have that $S_j = S_j'$ and the claim follows by the inductive hypothesis.
\end{proof}

\begin{claim}\label{claim:weak-to-strong:correctness:main:hyperedge-merge:property3}
     $\sum_{j \in [k]}\gamma_j\cdot \chi_{S_j} = \Tilde{w}_i(e^*, \cdot)$.
\end{claim}
\begin{proof}
    We consider two cases. First, suppose that $e^* \not = e\cup f$. Then, we have the following:
    $$\sum_{j \in [k]}\gamma_j\cdot \chi_{S_j} = \sum_{j \in [k']}\gamma_j'\cdot \chi_{S_j} = \Tilde{c}(e^*, \cdot) = \Tilde{w}(e^*, \cdot).$$
    Here, the first equality is because $\chi_{S_j} = \chi_{S_j'}$ since the hyperedges $e$ and $f$ are disjoint by \Cref{lem:WeakToStrongCover:properties-of-e-and-f-chosen-during-alg}. The second equality is by the inductive hypothesis. The final equality is by definition of the function $\Tilde{w}$. Next, suppose that the hyperedge $e^* = e\cup f$. We have the following:
    \begin{align*}
        \sum_{j \in [k]}\gamma_j\cdot \chi_{S_j}& = \sum_{j \in [k']}\gamma_j'\cdot \chi_{S_j'} + \alpha^R\chi_{e\cup f} &\\
        &= \sum_{j \in [k']}\gamma_j'\cdot \chi_{S_j'} + \alpha^R\chi_{S_k}&\\
        &= \Tilde{c}(e\cup f, \cdot) + \alpha^R\chi_{S_k}&\\
        &= \Tilde{w}(e\cup f, \cdot).&
    \end{align*}
    Here, the second equality is by definition of the set $S_i$. The third equality is by the inductive hypothesis. The final equality is by definition of the function $\Tilde{w}$.
\end{proof}

\begin{claim}\label{claim:weak-to-strong:correctness:main:hyperedge-merge:property4}
    $\sum_{j \in [k]}\gamma_j = w^*(e^*)$.
\end{claim}
\begin{proof}
    We consider two cases. First, suppose that $e^* \not = e\cup f$. Then, we have the following:
    $$\sum_{j \in [k]}\gamma_j = \sum_{j \in [k']}\gamma_j = c^*(e^*) = w^*(e^*).$$
    Here, the second equality is by the inductive hypothesis and the final equality is by \Cref{lem:WeakToStrong:properties-of-alpha^M-and-alpha^R}. Next, suppose that the hyperedge $e^* = e \cup f$. Then, we have the following:
    $$\sum_{j \in [k]}\gamma_j = \sum_{j \in [k']}\gamma_j' = c^*(e\cup f) + \alpha^R= w^*(e\cup f),$$
    where the second equality is by the inductive hypothesis.
\end{proof}
\fi

%% proof using recursion depth
\iffalse
\begin{proof}
    By \Cref{lem:weak-to-strong:correctness:main:strong-cover} the recursion depth of the algorithm is finite. Let $\ell \in \Z_+$ denote the recursion depth of the algorithm and let $i \in [\ell]$ be a recursive call. Let $(H_i = (V_i, E_i), w_i)$ and $p_i:2^{V_i}\rightarrow\Z$ denote the input hypergraph and input function to the $i^{th}$ recursive call. By \Cref{lem:weak-to-strong-cover:next-recursive-call-hypothesis-satification:symm-skew-supermod-strong-cover} and \Cref{lem:weak-to-strong-cover:next-recursive-call-hypothesis-satification:hyperedge-in-tight-set-positive-coverage} and induction on $i$, we have that the function $p_i$ is symmetric skew-supermodular, $b_{H_i, w_i}(X) \geq p_i(X)$ for every $X \subseteq V_i$, $\tau_{p_i}(e)\neq \emptyset$ for every hyperedge $e\in E_i$, and $b_{H_i, w_i}(u) > 0$ for every $u \in V_i$. Let $e_i, f_i\in E_i$ be the hyperedges chosen by the algorithm in the $i$th recursive call. Let $(H_i^*=(V, E_i^*), w_i^*)$ denote the hypergraph returned by the $i^{th}$ recursive call. We now show the claim by (reverse) induction on $i \in [\ell]$. For the base case, consider $i = \ell$. Here, we have that $(H_\ell^*, w_\ell^*) = (H_\ell, w_\ell) = (\emptyset, \emptyset)$ and so the claim holds. \knote{It would be better to do induction on  $\sum_{e\in E}w(e)$ instead of $i$. Such a proof will resemble the rest of the proofs in this section and would be easier to follow than figuring out this new index.} 
    
    For the inductive case, suppose that $i < \ell$.
    %By the inductive hypothesis, there exists a function $\Tilde{w}_{i+1}:E^*_{i+1}\times E_{i+1} \rightarrow \Z_{\geq 0}$, and for every $e^* \in E_{i+1}^*$ there exist positive integers $k_{i+1} \in \Z_+$ and $\gamma_1, \ldots, \gamma_{k_{i+1}} \in \Z_+$, and a collection of hyperedge subsets $S_1, \ldots S_{k_{i+1}} \subseteq E_{i+1}$ satisfying properties (1)-(4) of (\ref{tag:hyperedge-merge-properties}).
    For simplicity, we henceforth assume that  $e_i, f_i \in E_{i+1}$. We can make the previous assumption because if the hyperedge $e_i \not \in E_{i+1}$, then we can add the hyperedge into $E_{i+1}$ and set $w_{i+1}(e_i) \coloneqq  0$ \knote{Shouldn't the weight function be strictly positive?}. We similarly handle the case when the hyperedge $f_i \not \in E_{i+1}$. We also assume that the hyperedge $e_{i} \cup f_i \in E_{i+1}^*$ and $e_{i} \cup f_i \in E_{i}$. As in the explanation of the previous assumption, if $e_i\cup f_i \not \in E_{i+1^*}$, then we add the hyperedge $e_i \cup f_i$ to $E^*_{i+1}$ and set $w^*_{i+1}(e_i\cup f_i) \coloneqq  0$. We note that $e_{i}\cup f_i \not \in E_i$ and so we add the hyperedge to $E_i$ with $w_i(e_i\cup f_i) \coloneqq  0$. With these assumption and by \Cref{lem:WeakToStrong:properties-of-alpha^M-and-alpha^R}, we have that $E_i = E_{i+1}$. Furthermore, by the algorithm description, we have that  $E_{i}^* = E_{i+1}^*$.

    Let $e^* \in E_i^*$ be an arbitrary hyperedge.
    By the inductive hypothesis, there exists a function $\Tilde{w}_{i+1}:E^*_{i+1}\times E_{i+1} \rightarrow \Z_{\geq 0}$, positive integers $k' \in \Z_+$ and $\gamma_1', \ldots, \gamma_{k'}' \in \Z_+$, and a collection of hyperedge subsets $S_1', \ldots S_{k'}' \subseteq E_{i+1}$ satisfying properties (1)-(4) of (\ref{tag:hyperedge-merge-properties}).
    We now define the weight function $\Tilde{w}_i:E^*_{i}\times E_{i} \rightarrow \Z_{\geq 0}$ as follows. 
    %Let $e^* \in E_{i+1}^*$, and let $k', \gamma_1', \ldots, \gamma_{k'}' \in \Z_+$, and  $S_1', \ldots S_{k'}' \subseteq E_{i+1}$ be  values for $e^*$ satisfying (\ref{tag:hyperedge-merge-properties}). 
    For every $e \in E_i$ we let



$$\Tilde{w}_i(e^*, e) \coloneqq  \begin{cases}
        \Tilde{w}_{i+1}(e^*, e)&\text{ if } e \not\in \{e_i, f_i\}\\
        \Tilde{w}_{i+1}(e^*, e) + \Tilde{w}_{i+1}(e^*, e_i \cup f_i)&\text{ if } e \in \{e_i, f_i\} \text{ and } e^* \not = e_i\cup f_i\\
        \Tilde{w}_{i+1}(e^*, e) + \alpha_i^R& \text{ otherwise (i.e., if $e \in \{e_i, f_i\}$ and $e^* = e_i \cup f_i$).}
    \end{cases}$$

    \Cref{claim:weak-to-strong:correctness:main:hyperedge-merge:property1} below says that the weight function $\Tilde{w}_i$ satisfies property (1) of \Cref{def:WeakToStrongCover:hyperedge-merge}. For the hyperedge $e^*$, we now define positive integers $k \in \Z_+$ and $\gamma_1, \ldots, \gamma_{k} \in \Z_+$, and the collections of hyperedge subsets $S_1, \ldots S_{k} \subseteq E_{i}$ as follows:
    \begin{align*}
        k &: =\begin{cases}
            k' + 1& \text{ if } e^* = e_i\cup f_i\\
            k'&\text{ otherwise.}
        \end{cases}&\\
        \gamma_j&\coloneqq  \begin{cases}
            \alpha_i^R& \text{ if } e^* = e_i\cup f_i \text{ and } j = k \text{ for the hyperedge $e_i\cup f_i$}\\
            \gamma_j'&\text{ otherwise}.
        \end{cases}&\\
        S_j& \coloneqq  \begin{cases}
            \{e_i, f_i\}&\text{ if } e^* = e_i\cup f_i \text{ and } j = k \text{ for the hyperedge $e_i\cup f_i$}\\
            S_j' - \{e_i\cup f_i\} + \{e_i, f_i\}&\text{ if } e_i\cup f_i \in S_j'\\
            S_j'& \text{ otherwise.}
        \end{cases}&
    \end{align*}
    \Cref{claim:weak-to-strong:correctness:main:hyperedge-merge:property2}, \Cref{claim:weak-to-strong:correctness:main:hyperedge-merge:property3} and \Cref{claim:weak-to-strong:correctness:main:hyperedge-merge:property4} below show that the function $\Tilde{w}_i$, and the quantities $k$, $\gamma_1, \ldots, \gamma_k$ and $S_1, \ldots, S_k$ satisfy properties (2a)-(2c) of \Cref{def:WeakToStrongCover:hyperedge-merge} for the hyperedge $e^*$ and thus complete the proof of the claim.
\end{proof}

    \begin{claim}\label{claim:weak-to-strong:correctness:main:hyperedge-merge:property1}
        For every $e \in E_i$, we have that $\sum_{e^* \in E_i^*}\Tilde{w}_{i}(e^*, e) = w_i(e)$.
    \end{claim}
    \begin{proof}
        We note that $e_i \cup f_i \not \in E_i$ by \Cref{lem:WeakToStrongCover:properties-of-e-and-f-chosen-during-alg}. Let $e \in E_i$. First, suppose that $e \not \in \{e_i, f_i\}$. Then, we have the following:
        $$\sum_{e^* \in E_i^*}\Tilde{w}_{i}(e^*, e) = \sum_{e^* \in E_{i+1}^*}\Tilde{w}_{i+1}(e^*, e) = w_{i+1}(e) = w_i(e).$$
        Here, the first and final equalities are because $e \not \in \{e_i, f_i, e_i\cup f_i\}$ and he second equality is by the inductive hypothesis. Next, suppose that $e \in \{e_i, f_i\}$. Then, we have the following:
        \begin{align*}
            \sum_{e^* \in E_i^*}\Tilde{w}_i(e^*, e)&= \sum_{e^* \in E_{i+1}^*}\Tilde{w}_{i+1}(e^*, e) + \sum_{e^* \in E_{i+1}^*}\Tilde{w}_{i+1}(e^*, e_i\cup f_i) + \alpha_i^R&\\
            &= w_{i+1}(e) + w_{i+1}(e_i\cup f_i) + \alpha_i^R&\\
            & = w_i(e) - \alpha_i^M + (\alpha_i^M - \alpha_i^R) + \alpha_i^R&\\
            &= w_i(e).  &
        \end{align*}
        Here, the second equality is by the inductive hypothesis and the third equality is by \Cref{lem:WeakToStrong:properties-of-alpha^M-and-alpha^R}.
    \end{proof}

\begin{claim}\label{claim:weak-to-strong:correctness:main:hyperedge-merge:property2}
    For every $j \in [k]$, we have that $e^* = \biguplus_{e \in S_j}e$.
\end{claim}
\begin{proof}
    We consider three cases based on the definition of the collection $S_i$. First, suppose that $e^* = e_i \cup f_i$ and $j = k$ (for the hyperedge $e_i\cup f_i$). Then, we have that $S_j = \{e_i, f_i\}$. We recall that $e_i$ and $f_i$ are disjoint hyperedges by \Cref{lem:WeakToStrongCover:properties-of-e-and-f-chosen-during-alg} and so the claim holds. Next, suppose that $e_i \cup f_i \in S_j'$. Then, we have that $\uplus_{e \in S_j} e = \uplus_{e \in S_j'} e = e^*$, where the second equality is  because $e_i$ and $f_i$ are disjoint by \Cref{lem:WeakToStrongCover:properties-of-e-and-f-chosen-during-alg} , and the final equality is by the inductive hypothesis. Finally, suppose that $e_i\cup f_i \not \in S_j'$. Then, we have that $S_j = S_j'$ and the claim follows by the inductive hypothesis.
\end{proof}

\begin{claim}\label{claim:weak-to-strong:correctness:main:hyperedge-merge:property3}
     $\sum_{j \in [k]}\gamma_j\cdot \chi_{S_j} = \Tilde{w}_i(e^*, \cdot)$.
\end{claim}
\begin{proof}
    We consider two cases. First, suppose that $e^* \not = e_i\cup f_i$. Then, we have the following:
    $$\sum_{j \in [k]}\gamma_j\cdot \chi_{S_j} = \sum_{j \in [k']}\gamma_j'\cdot \chi_{S_j} = \Tilde{w}_{i+1}(e^*, \cdot) = \Tilde{w}_i(e^*, \cdot).$$
    Here, the first equality is because $\chi_{S_j} = \chi_{S_j'}$ since the hyperedges $e_i$ and $f_i$ are disjoint by \Cref{lem:WeakToStrongCover:properties-of-e-and-f-chosen-during-alg}. The second equality is by the inductive hypothesis. The final equality is by definition of the function $\Tilde{w}_i$. Next, suppose that $e^* = e_i\cup f_i$. We have the following:
    \begin{align*}
        \sum_{j \in [k]}\gamma_j\cdot \chi_{S_j}& = \sum_{j \in [k']}\gamma_j'\cdot \chi_{S_j'} + \alpha_i^R\chi_{e_i\cup f_i} &\\
        &= \sum_{j \in [k']}\gamma_j'\cdot \chi_{S_j'} + \alpha_i^R\chi_{S_k}&\\
        &= \Tilde{w}_{i+1}(e_i\cup f_i, \cdot) + \alpha_i^R\chi_{S_k}&\\
        &= \Tilde{w}_i(e_i\cup f_i, \cdot).&
    \end{align*}
    Here, the second equality is by definition of the set $S_i$. The third equality is by the inductive hypothesis. The final equality is by definition of the function $\Tilde{w}_i$.
\end{proof}

\begin{claim}\label{claim:weak-to-strong:correctness:main:hyperedge-merge:property4}
    $\sum_{j \in [k]}\gamma_j = w_i^*(e^*)$.
\end{claim}
\begin{proof}
    We consider two cases. First, suppose that $e^* \not = e_i\cup f_i$. Then, we have the following:
    $$\sum_{j \in [k]}\gamma_j = \sum_{j \in [k']}\gamma_j = w^*_{i+1}(e^*) = w_i^*(e^*).$$
    Here, the second equality is by the inductive hypothesis and the final equality is by \Cref{lem:WeakToStrong:properties-of-alpha^M-and-alpha^R}. Next, suppose that $e^* = e_i \cup f_i$. Then, we have the following:
    $$\sum_{j \in [k]}\gamma_j = \sum_{j \in [k']}\gamma_j' = w^*_{i+1}(e_i\cup f_i) + \alpha_i^R= w_i^*(e_i\cup f_i),$$
    where the second equality is by the inductive hypothesis.
\end{proof}
\fi

\subsection{Number of Recursive Calls and Hypergraph Support Size}\label{sec:WeakToStrongCover:NumberOfHyperedges}

We first set up some notation that will be used in the remainder of this section. 
By \Cref{lem:weak-to-strong:correctness:main:strong-cover}, the number of recursive calls made by \Cref{alg:WeakToStrongCover} is finite.
We will use $\ell$ to denote the depth of recursion of the algorithm. We will refer to the recursive call at depth $i$ as \emph{recursive call $i$} or the \emph{$i^{th}$ recursive call}. Let $i \in [\ell]$ be a recursive call. We let the hypergraph $(H_i = (V_i, E_i), w_i:E_i \rightarrow \Z_+)$ and the function $p_i : 2^{V_i} \rightarrow \Z$ denote the inputs to the $i^{th}$ recursive call. Furthermore, we let $e_i, f_i, \beta_i^M, \alpha_i^M, (H_i^M, w_i^M), \beta_i^R, \alpha_i^R, (H_i^R, w_i^R), p_i^R, (H_0^{i}, w_0^{i}), E_0^i, \zeros_i, V_i', p'_i$ denote the quantities $e, f, \beta^M, \alpha^M, (H^M, w^M), \beta^R, \alpha^R, (H^R, w^R), p^R, (H_0, w_0), E_0, \zeros, V', p'$ as defined by \Cref{alg:WeakToStrongCover} during the $i^{th}$ recursive call.
%Furthermore, we let $e_i, f_i \in E_i$ denote the hyperedges chosen by the algorithm, $\alpha_i^M, \alpha_i^R$ denote the values $\alpha^M, \alpha^R$, the set $\zeros_i$ denotes the set $\zeros$, $p_i^R \coloneqq  p_i - \alpha_i^R d_{(V, \{e_i\cup f_i\})}$, and $w_i^M, w_i^R$ denote the weight functions $w^M, w^R$ computed by the algorithm at the $i^{th}$ recursive call. 

%We recall that  a set $X\subseteq V$ is said to be \emph{tight} w.r.t. a function $p:2^V\rightarrow\Z$ and a hypergraph $(H = (V, E), w)$ if $p(X) = b_{(H,w)}(X)$. We also recall that $\tightSetFamily{p,H,w}$ denotes the family of tight sets and $\maximalTightSetFamily{p, H, w}$ denotes the family of \emph{inclusion-wise maximal} tight sets w.r.t. the function $p$ and hypergraph $(H,w)$. 
We now define certain set families of that will be crucial in our analysis. We let $\tightSetFamily{i}, \tightSetFamilyAfterMerge{i}, \tightSetFamilyAfterReduce{i}$ denote the families $\tightSetFamily{p_i, H_i, w_i}, \tightSetFamily{p_i, H_i^M, w_i^M}$ and $ \tightSetFamily{p_i^R, H_i^R, w_i^R}$ respectively, and $\maximalTightSetFamily{i}, \maximalTightSetFamilyAfterMerge{i}, \maximalTightSetFamilyAfterReduce{i}$ denote the families $\maximalTightSetFamily{p_i, H_i, w_i}, \maximalTightSetFamily{p_i, H_i^M, w_i^M}$ and $ \maximalTightSetFamily{p_i^R, H_i^R, w_i^R}$ respectively. We define the family of \emph{cumulative projected maximal $(p_i, H_i, w_i)$-tight sets} as:
$$\cumulativeMaximalTightSetFamily{i} \coloneqq  \begin{cases}
    \maximalTightSetFamily{1}& \text{ if $i = 1$,}\\
    
   \cumulativeMaximalTightSetFamily{i-1}|_{V_i} \cup \maximalTightSetFamily{i}& \text{ if $i \geq 2$}.
\end{cases}$$
\iffalse
$$\cumulativeMaximalTightSetFamily{i} \coloneqq  \begin{cases}
    \maximalTightSetFamily{1}& \text{ if $i = 1$,}\\
    
   \cumulativeMaximalTightSetFamily{i-1}|_{V_i} \cup \maximalTightSetFamily{i}& \text{ if $i \geq 2$}.
\end{cases}$$
\fi
The next lemma states two important properties of these set families.
\begin{lemma}\label{lem:WeakToStrong:properties-of-cumulative-tight-set-families}
Suppose that the input to \Cref{alg:WeakToStrongCover} is a non-empty hypergraph $(H = (V, E), w)$ and a symmetric skew-supermodular function $p:2^V\rightarrow \Z$ 
such that every hyperedge $\Tilde{e} \in E$ is contained in some $(p, H, w)$-tight set, $b_{(H,w)}(u) \not = 0$ for every $u \in V$, and $b_{(H,w)}(X) \geq p(X)$ for every $X \subseteq V$. Furthermore, suppose that the execution of \Cref{alg:WeakToStrongCover} on the input instance terminates using $\ell \in \Z_+$ recursive calls. Then, for every $i \in [\ell - 1]$, we have the following:
\iffalse
    Consider an execution of \Cref{alg:WeakToStrongCover} with $\ell \in \Z_+$ recursive calls. Suppose that the input to the $i^{th}$ recursive call is a hypergraph $(H_i = (V_i, E_i), w_i)$ and a function $p_i:2^{V_i}\rightarrow \Z$.
    Furthermore, suppose that  $(H_1, w_1)$ is a non-empty hypergraph and $p_1$ is a symmetric skew-supermodular function %satisfying conditions \ref{cond:WeakToStrong:(a)} and \ref{cond:WeakToStrong:(b)} 
such that every hyperedge $\Tilde{e} \in E_1$ is contained in some $(p_1, H_1, w_1)$-tight set, $b_{H_1, w_1}(u) \not = 0$ for every $u \in V_1$, and $b_{H_1,w_1}(X) \geq p_1(X)$ for every $X \subseteq V_1$.  Then, for every $i \in [\ell - 1]$, we have the following:
\fi
     \begin{enumerate}
        % \item $\tightSetFamily{i} \subseteq \tightSetFamilyAfterMerge{i} \subseteq \tightSetFamilyAfterReduce{i}$
        % \item $\projTightSetFamilyAfterReduce{i}{i+1} \subseteq \tightSetFamily{i+1}$
         \item $\projCumulativeMaximalTightSetFamily{i}{i+1} \subseteq \tightSetFamily{i+1}$, and
         \item the family $\cumulativeMaximalTightSetFamily{i}$ is laminar.
     \end{enumerate}
\end{lemma}
\begin{proof} 
Let $i\in [\ell - 1]$ be a recursive call. We recall that $p_i:2^{V_i}\rightarrow\Z$ and $(H_i, w_i)$ are inputs to the $i^{th}$ recursive call, where $p_1 := p$ and $(H_1, w_1) := (H, w)$.
By Lemmas \ref{lem:weak-to-strong-cover:next-recursive-call-hypothesis-satification:symm-skew-supermod-strong-cover} and \ref{lem:weak-to-strong-cover:next-recursive-call-hypothesis-satification:hyperedge-in-tight-set-positive-coverage} and induction on $i$, we have that $p_i$ is a symmetric skew-supermodular function and $(H_i = (V_i, E_i), w_i)$ is a non-empty hypergraph such that every hyperedge $\Tilde{e} \in E_i$ is contained in some $(p_i, H_i, w_i)$-tight set, $b_{(H_i, w_i)}(u) \not = 0$ for every $u \in V_i$, and $b_{(H_i, w_i)}(X) \geq p_i(X)$ for every $X \subseteq V_i$.
% We recall that the function $p_{i} = \functionContract{(p_{i-1} - b_{(H_0^{i-1}, w_0^{i-1})})}{\zeros_{i-1}}$ for every $i \geq 2$. Since $p_1$ is a symmetric skew-supermodular function, by \Cref{lem:weak-to-strong-cover:next-recursive-call-hypothesis-satification:symm-skew-supermod-strong-cover} and induction on $i$, we have that $p_{i}:2^{V_{i}}\rightarrow\Z$ is also a symmetric skew-supermodular function for every $i\ge 2$. 
% By \Cref{lem:WeakToStrongCover:TightSetFamiliesDuringAlgorithm} and induction on $i$, we have that $T_i\subseteq T_i^M \subseteq T_i^R$ and $T_i^R|_{V_{i+1}}\subseteq T_{i+1}$ for every $i\ge 2$. 
We now prove both claims separately below.
    \begin{enumerate}
        % \item We will first show that $\tightSetFamily{i}\subseteq \tightSetFamilyAfterReduce{i}$. Let $X \in \tightSetFamily{i}$. Here, we observe that $e_i\cap X, f_i\cap X \not = \emptyset$ as otherwise we have that 
        % $b_i^M(X) = b_i(X) - \alpha_i^M = p_i(X) - \alpha_i^M < p_i(X)$, where the first equality is because of the contradiction assumption and the final inequality is by \Cref{lem:WeakToStrong:properties-of-alpha^M-and-alpha^R}. This gives a contradiction. Thus, we have that $p_{i}(X) = b_{i}(X) = b_{i}^M(X)$, i.e. $X \in \tightSetFamilyAfterMerge{i}$.

        % We will now show the second inclusion in the claim, i.e. $\tightSetFamilyAfterMerge{i} \subseteq\tightSetFamilyAfterReduce{i}$. Let $X \in \tightSetFamilyAfterMerge{i}$. First, suppose that $\alpha_i^R = 0$. Then, we have that $b_i^R(X) = b_i^M(X) = p_i(X) = p_i^R(X)$, where the first and final equalities are because $\alpha_i^R = 0$. Thus, $X \in \tightSetFamilyAfterReduce{i}$. Next, suppose that $\alpha_{i}^R > 0$. Here, we consider the case where $e_i\cup f_i \not \in \delta_H(X)$ as otherwise, we have that $b_i^R(X) = b_i^M(X) - \alpha_i^R = p_i(X) - \alpha_i^R = p_i^R(X)$, and the claim holds. Then, we have the following: $b_{i}^R(X) = b_i^M(X) - \alpha_i^R = p_i^M(X) - \alpha_i^R < p_i^M(X) = p_i^R(X)$ which is a contradiction. Here, the final equality is because $d_{(V, \{e_i\cup f_i\})}(X) = 0$ since $e_i\cup f_i \not \in \delta_H(X)$.

        % \item Let $X \in \tightSetFamilyAfterReduce{i}$.  Then, we have the following:
        % $$p_{i+1}(X - \zeros_i) = p_i^R|_{\zeros_i}(X - \zeros_i) \geq p_i^R(X) = b_i^R(X) = b_i^R(X - \zeros_i),$$
        % where the final equality is because $\zeros_i = \{u \in V_i : b_i^R(u) = 0\}$. Thus, the set $X - \zeros_i \in \projTightSetFamilyAfterReduce{i}{i+1}$.
        \item By way of contradiction, let $i \in [\ell - 1]$ be the least index such that $\cumulativeMaximalTightSetFamily{i}|_{V_{i+1}} - T_{i+1} \not = \emptyset$. We note that $i \geq 2$ since $\cumulativeMaximalTightSetFamily{1} = \maximalTightSetFamily{1}\subseteq T_1$ and $T_1|_{V_2}\subseteq T_2$ by Lemma  \ref{lem:WeakToStrongCover:TightSetFamiliesDuringAlgorithm}(1) and \ref{lem:WeakToStrongCover:TightSetFamiliesDuringAlgorithm}(2), and hence, $\cumulativeMaximalTightSetFamily{1}|_{V_2}\subseteq T_2$. Thus, by our choice of index $i$, we have that $\cumulativeMaximalTightSetFamily{i-1}|_{V_i} \subseteq \tightSetFamily{i}$, and so it follows that $\cumulativeMaximalTightSetFamily{i} = \cumulativeMaximalTightSetFamily{i-1}|_{V_i} \cup \maximalTightSetFamily{i} \subseteq \tightSetFamily{i}$. Moreover, $T_i|_{V_{i+1}}\subseteq T_{i+1}$ by Lemma  \ref{lem:WeakToStrongCover:TightSetFamiliesDuringAlgorithm}(1) and \ref{lem:WeakToStrongCover:TightSetFamiliesDuringAlgorithm}(2). Consequently we have that $\cumulativeMaximalTightSetFamily{i}|_{V_{i+1}} \subseteq \tightSetFamily{i}|_{V_{i+1}} \subseteq \tightSetFamily{i+1}$, a contradiction to the choice of the index $i$. 
        %We also have that $\maximalTightSetFamily{i+1}\subseteq \tightSetFamily{i+1}$. Consequently, $\cumulativeMaximalTightSetFamily{i}|_{V_{i+1}}\cup \maximalTightSetFamily{i+1}\subseteq \tightSetFamily{i+1}$, a contradiction to the choice of the index $i$. 

        \iffalse
        We prove the claim by induction on $i$. 
        %We recall that $\cumulativeMaximalTightSetFamily{1} = \maximalTightSetFamily{1}$, and $\cumulativeMaximalTightSetFamily{i} = \left(\cumulativeMaximalTightSetFamily{i-1} \cup \maximalTightSetFamily{i-1}^M \cup \maximalTightSetFamily{i-1}^R\right)|_{V_i}\cup \maximalTightSetFamily{i}$ for every $i \geq 2$. 
        For the base case of $i=1$, we note that $\maximalTightSetFamily{1}\subseteq T_1$ and $T_1|_{V_2}\subseteq T_2$ by Lemma  \ref{lem:WeakToStrongCover:TightSetFamiliesDuringAlgorithm}(1) and \ref{lem:WeakToStrongCover:TightSetFamiliesDuringAlgorithm}(2), and hence, $\maximalTightSetFamily{1}|_{V_2}\subseteq T_2$. Let $i\ge 2$. 
        By the induction hypothesis, we have that $\cumulativeMaximalTightSetFamily{i-1}|_{V_i}\subseteq T_i$. Thus, $\cumulativeMaximalTightSetFamily{i} = \cumulativeMaximalTightSetFamily{i-1}|_{V_i} \cup \maximalTightSetFamily{i} \subseteq \tightSetFamily{i}$. Moreover, $T_i|_{V_{i+1}}\subseteq T_{i+1}$ by Lemma  \ref{lem:WeakToStrongCover:TightSetFamiliesDuringAlgorithm}(1) and \ref{lem:WeakToStrongCover:TightSetFamiliesDuringAlgorithm}(2). Consequently we have that $\cumulativeMaximalTightSetFamily{i}|_{V_{i+1}} \subseteq \tightSetFamily{i}|_{v_{i+1}} \subseteq \tightSetFamily{i+1}$.
        %Hence, $\cumulativeMaximalTightSetFamily{i-1}|_{V_i}|_{V_{i+1}}\subseteq T_{i+1}$. We also have that $\maximalTightSetFamily{i-1}^M\subseteq T_{i-1}^M\subseteq T_{i-1}^R$, $T_{i-1}^R|_{V_i}\subseteq T_i$, and $T_{i}|_{V_{i+1}}\subseteq T_{i+1}$, and hence $\maximalTightSetFamily{i-1}^M|_{V_i}|_{V_{i+1}}\subseteq T_{i+1}$. Similarly, we have that $\maximalTightSetFamily{i-1}^R\subseteq T_{i-1}^R$, $T_{i-1}^R|_{V_i}\subseteq T_i$, and $T_{i}|_{V_{i+1}}\subseteq T_{i+1}$, and hence $\maximalTightSetFamily{i-1}^R|_{V_i}|_{V_{i+1}}\subseteq T_{i+1}$. We have that $\maximalTightSetFamily{i}\subseteq T_i\subseteq T_i^M\subseteq T_i^R$ and $T_i^R|_{V_{i+1}}\subseteq T_{i+1}$, and hence, $\maximalTightSetFamily{i}|_{V_{i+1}}\subseteq T_{i+1}$. 
        \fi
        \item By way of contradiction, let $i \in [\ell - 1]$ be the least index such that the family $\cumulativeMaximalTightSetFamily{i}$  is not laminar. 
        %First, suppose that \cumulativeMaximalTightSetFamily{i} is not laminar.
        We note that $i \geq 2$ since the family $\cumulativeMaximalTightSetFamily{1} = \maximalTightSetFamily{1}$ is disjoint by \Cref{lem:weak-to-strong:uncrossing:maximal-tight-set-family-disjoint}. Thus, we have that $\cumulativeMaximalTightSetFamily{i} = \projCumulativeMaximalTightSetFamily{i-1}{i} \cup \maximalTightSetFamily{i}$.
        We note that the family $\maximalTightSetFamily{i}$ is also disjoint by \Cref{lem:weak-to-strong:uncrossing:maximal-tight-set-family-disjoint}. Furthermore, the family $\cumulativeMaximalTightSetFamily{i-1}$ (and consequently, the family  $\cumulativeMaximalTightSetFamily{i-1}|_{V_i}$) is laminar by our choice of $i\in[\ell - 1]$. Since the family $\cumulativeMaximalTightSetFamily{i}$ is not laminar, there exist distinct sets $X \in \projCumulativeMaximalTightSetFamily{i-1}{i}$ and $Y \in \maximalTightSetFamily{i}$ such that $X - Y, Y - X, X \cap Y \not = \emptyset$. By part (1) of the current lemma (shown above), we have that $X \in \tightSetFamily{i}$. Consequently, there exists $Z\in \maximalTightSetFamily{i}$ such that $X \subseteq Z$. We note that $Y \not = Z$ because $X -Z=\emptyset$ and $X - Y \not = \emptyset$. Thus, by \Cref{lem:weak-to-strong:uncrossing:maximal-tight-set-family-disjoint} we have that $Z\cap Y = \emptyset$. Therefore, $X\cap Y = \emptyset$, contradicting our choice of sets $X$ and $Y$.

        \iffalse
        \shubhang{Old proof with bug below.} By way of contradiction, let $i \in [\ell]$ be the least index such that $\cumulativeMaximalTightSetFamily{i}$ is not laminar. We note that $i \geq 2$ by \Cref{lem:weak-to-strong:uncrossing:maximal-tight-set-family-disjoint}. Then, there exist two sets $X, Y \in \cumulativeMaximalTightSetFamily{i}$ such that $X- Y, Y - X, X\cap Y \not = \emptyset$. We recall that $\cumulativeMaximalTightSetFamily{i} = \left(\cumulativeMaximalTightSetFamily{i-1} \cup \maximalTightSetFamily{i-1}^M \cup \maximalTightSetFamily{i-1}^R\right)|_{V_i}\cup \maximalTightSetFamily{i}$. By \Cref{lem:weak-to-strong:uncrossing:maximal-tight-set-family-disjoint}, both sets $X$ and $Y$ cannot be contained in the family $\maximalTightSetFamily{i}$. 

        First, suppose that $X, Y \in \left(\cumulativeMaximalTightSetFamily{i-1} \cup \maximalTightSetFamily{i-1}^M \cup \maximalTightSetFamily{i-1}^R\right)|_{V_i}$. Then, there exist sets $A_X, A_Y \in \cumulativeMaximalTightSetFamily{i-1} \cup \maximalTightSetFamily{i-1}^M \cup \maximalTightSetFamily{i-1}^R$ such that the sets $A_X-A_X, A_Y-A_X, A_X\cap A_Y\neq \emptyset$. If $A_X, A_Y \in \cumulativeMaximalTightSetFamily{{i-1}}$, then the index $i-1$ contradicts our choice of index $i$. Hence, at most one of the sets among $A_X$ and $A_Y$ can be in $\cumulativeMaximalTightSetFamily{{i-1}}$. 
        %Furthermore, if both sets $X, Y \in \projCumulativeMaximalTightSetFamily{i-1}{i}$, then there exist $A_X, A_Y \in \cumulativeMaximalTightSetFamily{i-1}$ such that $X \subseteq A_X$, $Y \subseteq A_Y$. Consequently the sets $A_X$ and $A_Y$ cross and the index $i-1$ contradicts our choice of index $i$. 
        We first consider the case where exactly one of the sets among $A_X$ and $A_Y$ is in $\cumulativeMaximalTightSetFamily{{i-1}}$. Without loss of generality, suppose that $A_X \in \cumulativeMaximalTightSetFamily{i-1}$. Suppose that $A_Y \in \maximalTightSetFamilyAfterMerge{i-1}$ (or $A_Y\in \maximalTightSetFamilyAfterReduce{i-1}$). Then, by \Cref{lem:WeakToStrongCover:TightSetFamiliesDuringAlgorithm}(1), we have that $\cumulativeMaximalTightSetFamily{i-1} \subseteq \tightSetFamily{i-1} \subseteq \tightSetFamilyAfterMerge{i-1}$ ($\subseteq \tightSetFamilyAfterReduce{i-1}$). In particular, there exists a set $B_X \in \maximalTightSetFamilyAfterMerge{i-1}$ (or $B_X \in \maximalTightSetFamilyAfterReduce{i-1}$) such that $A_X\subseteq B_X$. Since $A_X \cap A_Y \not = \emptyset$, we have that $B_X \cap A_Y \not = \emptyset$. Thus, by \Cref{lem:weak-to-strong:uncrossing:maximal-tight-set-family-disjoint}, we have that $A_X \subseteq B_X = A_Y$, contradicting $A_X-A_Y\neq \emptyset$. 
        Next, we consider the case where both $A_X$ and $A_Y$ are not in $\cumulativeMaximalTightSetFamily{{i-1}}$. By \Cref{lem:weak-to-strong:uncrossing:maximal-tight-set-family-disjoint}, both the sets $A_X$ and $A_Y$ cannot belong to the same family $\maximalTightSetFamilyAfterMerge{i-1}$ or $\maximalTightSetFamilyAfterReduce{i-1}$\shubhang{this is a bug. to use \Cref{lem:weak-to-strong:uncrossing:maximal-tight-set-family-disjoint}, every hyperedge needs to be contained in a tight set and every vertex needs to be contained in a hyperedge...}. Without loss of generality, let $A_X\in \maximalTightSetFamily{i-1}^M $ and $A_Y\in \maximalTightSetFamily{i-1}^R$. Then, by \Cref{lem:WeakToStrongCover:TightSetFamiliesDuringAlgorithm} part (1), we have that $\tightSetFamilyAfterMerge{i-1}\subseteq \tightSetFamilyAfterReduce{i-1}$. In particular, there exists a set $B_X \in \maximalTightSetFamilyAfterReduce{i-1}$ such that $A_X\subseteq B_X$. Since $A_X \cap A_Y \not = \emptyset$, we have that $B_X \cap A_Y \not = \emptyset$. Thus, by \Cref{lem:weak-to-strong:uncrossing:maximal-tight-set-family-disjoint}, we have that $A_X \subseteq B_X = A_Y$, contradicting $A_X-A_Y\neq \emptyset$. 
        %Using similar contradiction arguments, we can show that $A_X, A_Y \not \in \cumulativeMaximalTightSetFamily{i-1} \cup \maximalTightSetFamily{i-1}^M \cup \maximalTightSetFamily{i-1}^R$, a contradiction.
        
        
        Next, without loss of generality, suppose that $X \in \left(\cumulativeMaximalTightSetFamily{i-1} \cup \maximalTightSetFamily{i-1}^M \cup \maximalTightSetFamily{i-1}^R\right)|_{V_i}$ and $Y \in \maximalTightSetFamily{i}$. By \Cref{lem:WeakToStrongCover:TightSetFamiliesDuringAlgorithm} part (2), we have that $X \in \tightSetFamily{i}$. Thus, there exists a set $Z \in \maximalTightSetFamily{i}$ such that $X \subseteq Z$. Since $X-Y\neq \emptyset$, we have that $Z \not = Y$. Then, we have that $X\cap Y \subseteq Z \cap Y = \emptyset$, where the final equality is by \Cref{lem:weak-to-strong:uncrossing:maximal-tight-set-family-disjoint}. This contradicts $X\cap Y\neq \emptyset$.
        \fi
    \end{enumerate} 
\end{proof}

% \begin{enumerate}
%     \item $\tightSetFamily{i} \subseteq \tightSetFamilyAfterCreate{i}$,
%     \item  $\projTightSetFamilyAfterCreate{i}{i+1} \subseteq \tightSetFamily{i+1}$,  
%      \item $\cumulativeMaximalTightSetFamily{i} \subseteq \tightSetFamily{i}$,
%      \item $\cumulativeMaximalTightSetFamilyAfterCreate{i} \subseteq \tightSetFamilyAfterCreate{i}$, and
%      \item the families \cumulativeMaximalTightSetFamily{i} and \cumulativeMaximalTightSetFamilyAfterCreate{i} are laminar.
%  \end{enumerate}
% \begin{proof} 
% We recall that the function $p_{i} = (p_{i-1} - \alpha_{i-1} b_{(V_{i-1}, \{A_{i-1}\})})|_{\zeros_{i-1}}$ for $i \geq 2$. Since $p_1$ is a skew-supermodular function, we have that $p_{i}:2^{V_{i}}\rightarrow\Z$ is a skew-supermodular function for every $i\in [\ell]$. Furthermore, by \Cref{lem:Covering-Algorithm:hypothesis-after-contractionstep}, the functions $p_i$ and 
% $m_i:V_i\rightarrow\Z_+$ 
%     satisfy conditions (a) and (b) of \Cref{thm:WeakCoverViaUniformHypergraph:main} for every $i\in [\ell]$ since the functions $p_1$ and $m_1$ satisfy the two conditions. Let $i \in [\ell]$ be a recursive call. We now prove each claim separately below.
%     \begin{enumerate}
%         \item Let $X\in \tightSetFamily{i}$. We first show that $|A_i \cap X| \leq 1$. By way of contradiction, let $|A_i\cap X| \geq 2$. Then, we have the following:
%         \begin{align*}
%             m_{i}'(X) &= m_i(X) - \alpha_i\chi_{A_i}(X)&\\
%             &= p_i(X) - \alpha_i\chi_{A_i}(X)&\\
%             &\leq p_i(X) - 2\alpha_i&\\
%             &< p_i(X) - \alpha_i b_{(V_i, \{A_i\})}(X)&\\
%             & = p_{i}'(X),&
%         \end{align*}
%     which is a contradiction to \Cref{lem:Covering-Algorithm:hypothesis-after-createstep}. Here, the second equality is because $X \in \tightSetFamily{i}$. The first inequality is by our assumption that $|A_i\cap X|\ge 2$. The second inequality is because $b_{(V_i, \{A_i\})}(X) = 1$. We now show that $X \in \tightSetFamilyAfterCreate{i}$. We have the following: 
%     $$m_i'(X) = m_i(X) - \alpha_i|A_i\cap X| = p_i(X) - \alpha_i b_{(V_i, \{A_i\})}(X) = p_i'(X).$$
%     Here, the second equality holds because $X \in \tightSetFamily{i}$, and $b_{(V_i, \{A_i\})}(X) = |A_i\cap X|$ since $|A_i\cap X|\leq 1$. 

%     \item  We note that each set in the family $\projTightSetFamilyAfterCreate{i}{i+1}$ is of the form $Y - \zeros_i$ for some set $Y \in \tightSetFamilyAfterCreate{i}$. Thus, it suffices to consider an arbitrary set $X \in \tightSetFamilyAfterCreate{i}$ and show that $X - \zeros_i \in \tightSetFamily{i+1}$. We have the following:
%     $$m_{i+1}(X - \zeros_i) = m_i'(X) = p_i'(X) \leq p_{i+1}(X - \zeros_i) \leq m_{i+1}(X - \zeros_i).$$
%         Thus, all inequalities are equations and we have that $X - \zeros \in \tightSetFamily{i+1}$. Here, the first equality is because $m'(\zeros_i) = 0$. The second equality is because $X \in \tightSetFamilyAfterCreate{i}$. The first inequality is because $p_{i+1} = p_i'|_{\zeros_i}$. The final inequality is because $m_{i+1}(Z) \geq p_{i+1}(Z)$ for every $Z\subseteq V_{i+1}$ by \Cref{lem:Covering-Algorithm:hypothesis-after-contractionstep}. We remark that the statement (and proof) of this part can be strengthened to say that $\projTightSetFamilyAfterCreate{i}{i+1} = \tightSetFamily{i+1}$, but the weaker statement suffices for our purposes and so we omit the details. 
        
%         \item By way of contradiction, let $i \in [\ell]$ be the smallest recursive call such that the claim is false. We note that $i \geq 2$ because $\cumulativeMaximalTightSetFamily{1} = \maximalTightSetFamily{1} \subseteq \tightSetFamily{1}$. Thus, we have that $\cumulativeMaximalTightSetFamily{i} = \projCumulativeMaximalTightSetFamily{i-1}{i} \cup \projMaximalTightSetFamilyAfterCreate{i-1}{i} \cup \maximalTightSetFamily{i}$. We observe that  $\cumulativeMaximalTightSetFamily{i-1}\subseteq \tightSetFamily{i-1} \subseteq \tightSetFamilyAfterCreate{i-1}$, where the first inclusion is by our choice of $i \in [\ell]$, and the second inclusion is by part (a). Furthermore, we have that  $\projCumulativeMaximalTightSetFamily{i-1}{i} \cup \projMaximalTightSetFamilyAfterCreate{i-1}{i} \subseteq \projTightSetFamilyAfterCreate{i-1}{i} \subseteq \tightSetFamily{i}$, where the second inclusion is by parts (a) and (b). Thus, $\cumulativeMaximalTightSetFamily{i} \subseteq \tightSetFamily{i}$.

%         \item We recall that $\cumulativeMaximalTightSetFamilyAfterCreate{i} = \cumulativeMaximalTightSetFamily{i} \cup \maximalTightSetFamilyAfterCreate{i}$. Then, we have that $\cumulativeMaximalTightSetFamily{i} \subseteq \tightSetFamilyAfterCreate{i}$ by parts (c) and (a), and $\maximalTightSetFamilyAfterCreate{i} \subseteq \tightSetFamilyAfterCreate{i}$ by definition.
        
%         \item By way of contradiction, let $i \in [\ell]$ be the smallest recursive call such that the family $\cumulativeMaximalTightSetFamily{i}$ or \cumulativeMaximalTightSetFamilyAfterCreate{i} is not laminar. First, suppose that \cumulativeMaximalTightSetFamily{i} is not laminar.
%         We note that $i \geq 2$ as the family $\cumulativeMaximalTightSetFamily{1} = \maximalTightSetFamily{1}$ is disjoint by \Cref{lem:PropertiesOfTightSets:calTpmDisjoint}. Thus, we have that $\cumulativeMaximalTightSetFamily{i} = \projCumulativeMaximalTightSetFamilyAfterCreate{i-1}{i} \cup \maximalTightSetFamily{i}$.
%         We note that the family $\maximalTightSetFamily{i}$ is disjoint by \Cref{lem:PropertiesOfTightSets:calTpmDisjoint}. Furthermore, the family $\cumulativeMaximalTightSetFamilyAfterCreate{i-1}$ is laminar by our choice of $i\in[\ell]$. Since the family $\cumulativeMaximalTightSetFamily{i}$ is not laminar, there exist distinct sets $X \in \cumulativeMaximalTightSetFamily{i-1}$ and $Y \in \maximalTightSetFamily{i}$ such that the sets $X$ and $Y$ cross. By part (c), we have that $X \in \tightSetFamily{i}$. Consequently, there exists $Z\in \maximalTightSetFamily{i}$ such that $X \subseteq Z$. By \Cref{lem:PropertiesOfTightSets:calTpmDisjoint} we have that $Z\cap Y = \emptyset$. Thus, $X\cap Y = \emptyset$, contradicting our choice of sets $X$ and $Y$.

%         Next, suppose that \cumulativeMaximalTightSetFamily{i} is laminar, but \cumulativeMaximalTightSetFamilyAfterCreate{i} is not laminar. We recall that $\cumulativeMaximalTightSetFamilyAfterCreate{i} = \cumulativeMaximalTightSetFamily{i} \cup \maximalTightSetFamilyAfterCreate{i}$. Similar to the previous case, there exist distinct sets  $X \in \cumulativeMaximalTightSetFamily{i}$ and $Y \in \maximalTightSetFamilyAfterCreate{i}$ such that the sets $X$ and $Y$ cross. By parts (c) and (a), we have that $X \in \tightSetFamilyAfterCreate{i}$. Thus, there exists a set $Z \in \maximalTightSetFamilyAfterCreate{i}$ such that $X \subseteq Z$. By \Cref{lem:PropertiesOfTightSets:calTpmDisjoint}, either $Y = Z$ or $Y\cap Z = \emptyset$. In particular, we have that either $X \subseteq Y$ or $X\cap Y = \emptyset$, contradicting our choice of the sets $X,Y$.
%     \end{enumerate}
% \end{proof}

\begin{lemma}\label{lem:WeakToStrongCover:NumberOfHyperedges}
%Consider an execution of \Cref{alg:WeakToStrongCover}.
Suppose that the input to \Cref{alg:WeakToStrongCover} is a non-empty hypergraph $(H = (V, E), w)$ and a symmetric skew-supermodular function $p:2^V\rightarrow \Z$ 
%satisfying conditions \ref{cond:WeakToStrong:(a)} and \ref{cond:WeakToStrong:(b)} 
such that every hyperedge $\Tilde{e} \in E$ is contained in some $(p, H, w)$-tight set, $b_{(H,w)}(u) \not = 0$ for every $u \in V$, and $b_{(H,w)}(X) \geq p(X)$ for every $X \subseteq V$. 
Then, we have the following:
    \begin{enumerate}
        \item the recursion depth $\ell \in \Z_+$ of the algorithm is at most $|E| + 10|V| - 1$, and
        \item $|E^*| \leq |E| + 10|V| - 2$, where $(H^* = (V, E^*), w^*)$ denotes the hypergraph returned by the algorithm.
        %\knote{This lemma needs to be strengthened to say that the number of hyperedges in the hypergraph input to every intermediate recursive call is at most $|E|+10|V|$. This fact is the fact used in the run-time analysis. }
    \end{enumerate}
\end{lemma}

\begin{proof}
Let $i \in [\ell - 1]$ be a recursive call. We recall that $(H_i = (V_i, E_i), w_i:E_i \rightarrow \Z_+)$ and $p_i : 2^{V_i} \rightarrow \Z$ denote the inputs to the $i^{th}$ recursive call, and $(H_i^* = (V, E_i^*), w_i^*)$ denotes the hypergraph returned by the $i^{th}$ recursive call. Furthermore, $e_i, f_i, \beta_i^M, \alpha_i^M, (H_i^M, w_i^M), \beta_i^R, \alpha_i^R, (H_i^R, w_i^R)$, $p_i^R$, $(H_0^{i}, w_0^{i})$, $E_0^i, \zeros_i$ denote the quantities  as defined by \Cref{alg:WeakToStrongCover} during the $i^{th}$ recursive call.
    We prove the claims separately below.
    \begin{enumerate}
        \item We have the following:
        $$\ell \leq \left|\left\{i \in [\ell - 1] : \alpha_i^R = \beta_i^R\right\}\right|  +  \left|\left\{i \in [\ell- 1] : \alpha_i^R = w_i^M(e_i \cup f_i)\right\}\right| + 1 \leq |E| + 10|V| - 1,$$
        where the final inequality is by \Cref{claim:WeakToStrong:hyperedge-support-size:alpha^R=weight-of-added-hyperedge} and \Cref{claim:WeakToStrong:hyperedge-support-size:alpha^R<weight-of-added-hyperedge} below.
        \item 
        %\shubhang{todo: the proof is faulty/incomplete as it does not account for hyperedges that were reduced and added back.} Consider a recursive call $i \in [\ell]$. We recall that by part (4) of \Cref{lem:WeakToStrong:properties-of-alpha^M-and-alpha^R}, we have that $E_{i+1} = E_i - \{e_i\}\cdot \indicator_{\alpha_i^M = w_i(e_i)} - \{f_i\}\cdot \indicator_{\alpha_i^M = w_i(f_i)} + \{e_i\cup f_i\}\cdot \indicator_{\alpha_i^R < w_i^M(e_i\cup f_i)}$. Consequently, regardless of the value of $\alpha_i^M$, we have that $|E_{i+1}| \leq |E_i| + 1$. Furthermore, the inequality is an equality if and only if $\alpha_i^M = w_i(e_i) = w_i (f_i)$ and $\alpha_i^R = \beta_i^R < w_i^M(e_i\cup f_i)$. The number of such recursive calls is at most $5|V| - 1$ by \Cref{claim:WeakToStrong:hyperedge-support-size:alpha^R=weight-of-added-hyperedge}.
        %Let $i \in [\ell]$ be a recursive call. 
        %Let $\left(H_i = (V_i, E_i), w_i\right)$ denote the hypergraph input to the $i^{th}$ recursive call, and $\left(H_i^* = (V_i, E_i^*), w_i^*\right)$ denote the hypergraph returned by the $i^{th}$ recursive call. We recall that by the algorithm, we have that $(H_i^*, w_i^*) = (H_{i+1}^*, w_{i+1}^*) + \left((V_i,\{e_i\cup f_i\}), \alpha^R_i\chi_{e_i\cup f_i}\right)$. Consequently, 
        We have the following:
        $$|E_1^*| \leq |E_2^*| + |E_0^2| =   |E_{2}^*| + \indicator_{\alpha_1^R > 0} = |E_{\ell}^*| + \sum_{i \in [\ell - 1]}\indicator_{\alpha_i^R > 0} \leq \ell - 1\leq |E_1| + 10|V_1| - 2.$$
        Here, the second inequality is because $E_{\ell}^* = \emptyset$ by definition of the \Cref{alg:WeakToStrongCover} and the final inequality is by part (1) of the claim (shown above).
    \end{enumerate}
\end{proof}

\begin{claim}\label{claim:WeakToStrong:hyperedge-support-size:alpha^R=weight-of-added-hyperedge}
$\left|\left\{i \in [\ell-1] : \alpha_i^R = \beta_i^R\right\}\right| \leq 5|V| - 1$.
\end{claim}
\begin{proof}
We define a potential function $\phi:[\ell - 1] \rightarrow\Z_{\geq 0}$ as follows: for every $i \in [\ell]$,
    $$\phi(i) \coloneqq  \left|\cumulativeMaximalTightSetFamily{i}\right| + 3|\cumulativeZeros{i-1}|,$$
    where we define the set $\cumulativeZeros{0} \coloneqq  \emptyset$.
We first show that the function $\phi$ is monotone non-decreasing.
Let $i \in [\ell -1]$ be a recursive call of \Cref{alg:WeakToStrongCover}. If $\zeros_i = \emptyset$, then $|\cumulativeMaximalTightSetFamily{i}| \leq |\cumulativeMaximalTightSetFamily{i+1}|$ since $\cumulativeMaximalTightSetFamily{i+1} = \cumulativeMaximalTightSetFamily{i} |_{V_{i+1}}\cup \maximalTightSetFamily{i+1} = \cumulativeMaximalTightSetFamily{i}  \cup \maximalTightSetFamily{i+1},$   where the second equality is because $V_{i+1} = V_i - \zeros_i = V_i$. Then, we have that $\phi(i+1) \geq \phi(i)$ since $\cumulativeZeros{i-1} \subseteq \cumulativeZeros{i}$ by definition. Thus, we suppose that $\zeros_i \not = \emptyset$. We have the following: 
\begin{align*}
    \phi(i+1) - \phi(i)& = |\cumulativeMaximalTightSetFamily{i+1}| + 3 |\cumulativeZeros{i}| - |\cumulativeMaximalTightSetFamily{i}| - 3|\cumulativeZeros{i-1}|&\\
    & = \left|\cumulativeMaximalTightSetFamily{i}|_{V_{i+1}}\cup \maximalTightSetFamily{i+1}  \right| - |\cumulativeMaximalTightSetFamily{i}| + 3|\zeros_{i}|&\\
    &\geq \left|\cumulativeMaximalTightSetFamily{i}|_{V_{i+1}}\right| - |\cumulativeMaximalTightSetFamily{i}| + 3|\zeros_{i}|&\\
    & \geq  0.&
\end{align*}
Here, the final inequality is because of the following: we recall that the family $\cumulativeMaximalTightSetFamily{i}$ is laminar by \Cref{lem:WeakToStrong:properties-of-cumulative-tight-set-families}. Furthermore, the family $\projCumulativeMaximalTightSetFamily{i}{i+1}$ is the projection of the laminar family $\cumulativeMaximalTightSetFamily{i}$ to the ground set $V_{i+1} = V_i - \zeros_i$. The inequality then follows by \Cref{lem:CoveringAlgorithm:projection-laminar-family}  --  here, we observe that  \Cref{lem:CoveringAlgorithm:projection-laminar-family} says that this inequality is strict because $\zeros_i \not = \emptyset$. We will refer back to this observation in the next part of the proof.

Consider a recursive call $i \in \left\{i \in [\ell - 1] : \alpha_i^R = \beta_i^R\right\}$. We now show that $\phi(i+1) > \phi(i)$, i.e. the potential value strictly increases. 
We suppose that $\zeros_i = \emptyset$ as otherwise $\phi(i+1) > \phi(i)$ by the observation in the final line of the previous paragraph.
Thus, $p_{i+1} = p_i^R$ and $(H', w') = (H^R, w^R)$. 
We note that the hypergraph $(H_i, w_i)$ is non-empty since $i < \ell$.
We consider a set $X\subseteq V_i$ such that $e_i\cup f_i\subseteq X\subseteq V_i$ and $b_i^M(X)-p_i(X)=\beta_i^R$ (such a set exists by the definition of $\beta_i^R$). 
Then, we have the following: $b_{i+1}(X) = b_i^R(X)=b_i^M(X)-\alpha_i^R=p_i(X)=p_i^R(X) = p_{i+1}(X)$, where the first equality is because $(H', w') = (H^R, w^R)$, the second equality is because $\zeros_i = \emptyset$, the third equality is because $\alpha_i = \beta_i$ by our choice of $i \in [\ell - 1]$, and the final equality is because $p_{i+1} = p_{i}^R$. Thus, $X \in \tightSetFamily{i+1}$. Consequently, 
%Since $\alpha_i^R=\beta_i^R$, we have that $b_i^R(X)=b_i^M(X)-\alpha_i^R=p_i(X)=p_i^R(X)$ and hence, $X\in \tightSetFamily{i}^R$. In particular, 
there exists a set $Y \in \maximalTightSetFamily{i+1}$ such that $e_i\cup f_i \subseteq X \subseteq Y$. We note that $Y \not \in \tightSetFamily{i}$ due to the following: suppose (for sake of contradiction) that $Y \in \tightSetFamily{i}$, then 
%In particular, we have that the set $\maximalTightSetContainingHyperedge{i}^R(e_i\cup f_i)\not = \emptyset$ because $e_i\cup f_i \subseteq X$. We note that the set $\maximalTightSetContainingHyperedge{i}^R(e_i\cup f_i) \not \in \tightSetFamily{i}$ due to the following: suppose that the set $Y\coloneqq \maximalTightSetContainingHyperedge{i}^R(e_i\cup f_i) \in \tightSetFamily{i}$, then
$b_i(Y)=p_i(Y)$ and $b_i^M(Y)=b_i(Y)-\alpha_i^M=p_i(Y)-\alpha_i^M<p_i(Y)$, a contradiction to \Cref{lem:WeakToStrong:properties-of-alpha^M-and-alpha^R}(3) -- the first equality is by definition of $b_i^M$ and $e_i\cup f_i\subseteq Y$ and the inequality is because $\alpha_i^M\ge 1$ by \Cref{lem:WeakToStrong:properties-of-alpha^M-and-alpha^R}(1). Thus, $Y \in \maximalTightSetFamily{i+1} - \tightSetFamily{i} \subseteq \maximalTightSetFamily{i+1} - \cumulativeMaximalTightSetFamily{i}$ by \Cref{lem:WeakToStrong:properties-of-cumulative-tight-set-families}(1).
Consequently, we have that
$\cumulativeMaximalTightSetFamily{i} \subsetneq \cumulativeMaximalTightSetFamily{i} \cup \maximalTightSetFamily{i+1} = \cumulativeMaximalTightSetFamily{i+1}.$
%Here, the strict inclusion is because $ \maximalTightSetFamily{i}^R - \cumulativeMaximalTightSetFamily{i} \not = \emptyset$ and the equality is because $V_{i+1} = V_i - \zeros = V_i$.
%since $\cumulativeMaximalTightSetFamily{i} \subseteq \tightSetFamily{i}$ by \Cref{lem:WeakToStrong:properties-of-cumulative-tight-set-families}. The first equality is because $\alpha_i^R = 0$ by \Cref{lem:WeakToStrong:properties-of-alpha^M-and-alpha^R} and $V_{i+1} = V_i - \zeros_i = V_i$. Thus, $\phi(i+1) > \phi(i)$.

We now show the claim. By the previous two properties of our potential function $\phi$, we have the following:
$$\left|\left\{i \in [\ell - 1] : \alpha_i^R = \beta_i^R\right\}\right| \leq \phi(\ell - 1) - \phi(1) \leq |\cumulativeMaximalTightSetFamily{\ell - 1}| + 3|\cumulativeZeros{\ell - 1}| \leq 5|V| - 1,$$
where the final inequality is because the family $\cumulativeMaximalTightSetFamily{\ell - 1}$ is laminar by \Cref{lem:WeakToStrong:properties-of-cumulative-tight-set-families}(2).
\end{proof}

% \begin{proof}
%     Let $\ell \in \Z_+$ denote the number of recursive calls of the execution of \Cref{alg:WeakToStrongCover}. Let $i \in [\ell]$ be a recursive call and let the function $p_i:2^{V_i}\rightarrow \Z$ and the hypergraph $(H_i = (V_i, E_i), w_i)$ denote the inputs to the $i^{th}$ recursive call. 
%     We define a potential function $\phi:[\ell] \rightarrow\Z_{\geq 0}$ as follows: for every $i \in [\ell]$,
%     $$\phi(i) \coloneqq  \left|\cumulativeMaximalTightSetFamily{i}\right| + 3|\cumulativeZeros{i-1}|,$$
%     where we define the set $\cumulativeZeros{0} \coloneqq  \emptyset$.
%      Let $(H^* = (V_i, E_i^*), w_i^*)$ denote the hypergraph returned by the $i^{th}$ recursive call. Then, \Cref{claim:WeakToStrong:hyperedge-support-szie:main} below shows that the function $\phi$ is monotone non-decreasing and that $\phi(i+1) > \phi(i)$ if $|E^*_{i}| > |E^*_{i+1}|$. We note that by the definition of \Cref{alg:WeakToStrongCover}, $|E^*_{i}| - |E^*_{i+1}| \leq 1$.
%     Thus, we have the following:
%     $$|E^*| - |E| \leq \sum_{i \in [\ell - 1]} \left(\phi(i+1) - \phi(i)\right)
%          \leq |\cumulativeMaximalTightSetFamily{\ell}| + 3|\cumulativeZeros{\ell - 1}|\leq 5|V| - 1,$$
% where the second inequality is because the function $\phi$ is non-decreasing and the final inequality is because $|\cumulativeMaximalTightSetFamily{\ell}| \leq 2|V| - 1$ since the family $\cumulativeMaximalTightSetFamily{\ell}$ is laminar by \Cref{lem:WeakToStrong:properties-of-cumulative-tight-set-families}.
% \end{proof}

\begin{claim}\label{claim:WeakToStrong:hyperedge-support-size:alpha^R<weight-of-added-hyperedge}
    $\left|\left\{i \in [\ell - 1] : \alpha_i^R = w_i^M(e_i \cup f_i)\right\}\right| \leq |E| + 5|V| - 1$. %\kristof{Maybe dismiss the index of $E$ and $V$?}
\end{claim}
\begin{proof}
    We define a potential function $\phi:[\ell - 1] \rightarrow \Z_{\geq 0}$ as follows: for every $i \in [\ell]$,
    $$\phi(i) \coloneqq  -|E_i| - \left|\left\{j \in [i,\ell - 1] : \alpha_j^R < w_j^M(e_j\cup f_j)\right\}\right|.$$
    
    First, consider a recursive call $i \in [\ell - 1]$ such that $\alpha_i^R = \beta_i^R < w_i^M(e_i \cup f_i)$. We will show that the potential function $\phi$ does not decrease. By Lemma \ref{lem:WeakToStrong:properties-of-alpha^M-and-alpha^R}(2) and \ref{lem:WeakToStrong:properties-of-alpha^M-and-alpha^R}(5), we have that 
    \begin{align*}
       E_{i+1} & = E_i - \{e_i\}\cdot \indicator_{\alpha_i^M = w_i(e_i)} - \{f_i\}\cdot \indicator_{\alpha_i^M = w_i(f_i)} + \{e_i\cup f_i\}\cdot \indicator_{\alpha_i^R = \beta_i^R < w_i^M(e_i \cup f_i)}&\\
        &= E_i - \{e_i\}\cdot \indicator_{\alpha_i^M = w_i(e_i)} - \{f_i\}\cdot \indicator_{\alpha_i^M = w_i(f_i)} + \{e_i\cup f_i\}.&
    \end{align*}
Consequently, 
    %regardless of the value of $\alpha_i^M$, 
    we have that $|E_{i+1}| \leq |E_i| + 1$. Furthermore, by the choice of recursive call $i$, we have that $\left|\{j \in [i+1, \ell - 1] : \alpha_j^R < w_j^M(e_j\cup f_j)\}\right| = \left|\{j \in [i, \ell - 1] : \alpha_j^R < w_j^M(e_j\cup f_j)\}\right| - 1$. Thus, $\phi(i+1) \geq \phi(i)$.
    
     Next, consider a recursive call $i \in [\ell - 1]$ such that $\alpha_i^R = w_i^M(e_i \cup f_i)$. We will show that the potential function $\phi$ strictly increases. We note that the hypergraph $(H_i, w_i)$ is non-empty since $i < \ell$.  
     By \Cref{lem:WeakToStrong:properties-of-alpha^M-and-alpha^R}(1), we have that $\alpha_i^R = w_i^M(e_i \cup f_i) = \alpha_i^M \geq 1$. 
     Consequently, by  \Cref{lem:WeakToStrong:alpha^R=0-iff-alpha^M=beta^M}, we have that $\alpha_i^M  < \beta_i^M$ and hence, $\alpha_i^M \in \{w_i(e_i), w_i(f_i)\}$ by definition. 
     Suppose that $\alpha_i^M = w_i(e_i)$. 
     Then, we have that $w^R_i(e_i) = w^M_i(e_i) = w_i(e_i) - \alpha_{i}^M = 0$. 
     Thus, we have that $e_i \not \in E_{i+1}$. 
     Then, by our choice of $i$ and by Lemma \ref{lem:WeakToStrong:properties-of-alpha^M-and-alpha^R}(2) and \ref{lem:WeakToStrong:properties-of-alpha^M-and-alpha^R}(5), we have that $E_{i+1} \subsetneq E_i$. 
     We note that this strict inclusion can also be obtained in the case where $\alpha_i^M = w_i(f_i)$ using the same argument. Furthermore, by the choice of recursive call $i$, we have that $\left|\{j \in [i+1, \ell - 1] : \alpha_j^R < w_j^M(e_j\cup f_j)\}\right| = \left|\{j \in [i, \ell - 1] : \alpha_j^R < w_j^M(e_j\cup f_j)\}\right|$.
     Thus, $\phi(i+1) > \phi(i)$.

      The previous two properties of our potential function give us the following:
      $$\left|\left\{i \in [\ell - 1] : \alpha_i^R = w_i^M(e_i \cup f_i)\right\}\right| \leq \phi(\ell - 1) - \phi(1) \leq |E_1| + \left|\left\{j \in [\ell - 1] : \alpha_j^R < w_j^M(e_j\cup f_j)\right\}\right| \leq |E| + 5|V| - 1,$$ 
      %\kristof{Maybe dismiss the index of $E$ and $V$?}
      where the final inequality is by \Cref{claim:WeakToStrong:hyperedge-support-size:alpha^R=weight-of-added-hyperedge} -- we note that if $\alpha_j^R < w_j^M(e_j\cup f_j)$, then $\alpha_j^R = \beta_j^R$ by definition.
\end{proof}

\iffalse
\subsection{Strongly Polynomial Runtime of \Cref{alg:WeakToStrongCover}}\label{sec:WeakToStrongCover:strongly-poly-runtime:main}
\shubhang{Todo.}
In this section, we show that \Cref{alg:WeakToStrongCover} can be implemented to run in strongly polynomial time, given access to \functionMaximizationOracleStrongCover{p}. As a necessary step to show this, the next lemma says that maximization oracle for the input function $p$ during a recursive call of \Cref{alg:WeakToStrongCover} can be used to implement the maximization oracle for the input function $p'$ to the subsequent recursive call in strongly polynomial time. We note that the proof of the lemma is very similar to the proof of \Cref{lem:p''-max-oracle-from-p-max-oracle} in \Cref{sec:CoveringAlgorithm:p''-max-oracle-using-p-max-oracle} and so we omit its proof for brevity.

\begin{lemma}\label{lem:WeakToStrong:p'-oracle-from-p-oracle}
Consider a recursive call of \Cref{alg:WeakToStrongCover}. Suppose that the input to the recursive call is a non-empty hypergraph $(H = (V, E), w)$ and a symmetric skew-supermodular function $p:2^V\rightarrow \Z$ satisfying conditions \ref{cond:WeakToStrong:(a)} and \ref{cond:WeakToStrong:(b)} such that $b_{(H,w)}(X) \geq p(X)$ for every $X \subseteq V$. Let $p':2^{V'}\rightarrow\Z$ be the function constructed by the algorithm as input to the subsequent recursive call. 
Then, a  \functionMaximizationOracleStrongCover{p'} can be implemented using a  \functionMaximizationOracleStrongCover{p} in strongly polynomial time.
\end{lemma}

The next lemma shows that the family $\maximalTightSetFamily{p, H, w}$ can be computed in strongly polynomial number of calls to \functionMaximizationOracleStrongCover{p}. We will subsequently use this in \Cref{lem:WeakToStrongCover:single-recursive-call-strongly-polytime} to show that a recursive call of \Cref{alg:WeakToStrongCover} runs in strongly polynomial time given the maximization oracle of the input function.

 \begin{lemma}\label{lem:WeakToStrong::maximal-tight-set-family-computed-in-strongly-polytime}
 Let $p:2^V\rightarrow\Z$ be a skew-supermodular function and let $\left(H = (V, E), w:E \rightarrow\Z_+\right)$ be a hypergraph satisfying conditions \ref{cond:WeakToStrong:(a)} and \ref{cond:WeakToStrong:(b)} such that $b_{(H,w)}(X) \geq p(X)$ for every $X \subseteq V$. Then, the family $\maximalTightSetFamily{p, H, w}$ can be computed in $\poly(|V|)$ calls to \functionMaximizationOracleStrongCover{p}.
\end{lemma}
\begin{proof}
    We note that we can find a $(p, H, w)$-tight set $X\subseteq V$ 
    %containing the set $Z \subseteq V$ 
    (if one exists) by a single call to \functionMaximizationOracleStrongCover{p} by letting  $(G_0, w_0) \coloneqq  \left((V, \emptyset), \emptyset\right)$ to be the empty hypergraph, the sets $S_0, T_0 \coloneqq  \emptyset$, and the vector $y_0$ to be the all-zeros vector as inputs to $\functionMaximizationOracleStrongCover{{p}}$. 
    Let $X \subseteq V$ be a $(p, H, w)$-tight set. Then, we can find a $(p, H, w)$-tight set $Y \subseteq V$ such that $X \subseteq Y$ (if such a set $Y$ exists) by guessing a vertex $u \in Y - X$ (by enumerating over all vertices $u \in V - X$) and querying $\functionMaximizationOracleStrongCover{p}$ on inputs  $(G_0, w_0) \coloneqq  \left((V, \emptyset), \emptyset\right)$, $S_0 \coloneqq  X \cup \{u\}$, $T_0 \coloneqq  \emptyset$, and the vector $y_0$ as the all-zeros vector. Iteratively applying this procedure will result in a maximal $(p, H, w)$-tight set $Z \in \maximalTightSetFamily{p, H, w}$. By \Cref{lem:weak-to-strong:uncrossing:maximal-tight-set-family-disjoint}, The family $\maximalTightSetFamily{p, H, w}$ is disjoint. Thus, recursively applying the procedure on the set $V - Z$ (i.e. including $Z$ in the parameter $T_0$ to all subsequent queries to \functionMaximizationOracleStrongCover{p}) will give us the family $\maximalTightSetFamily{p, H, w}$. This procedure requires  $\poly(|V|)$ many calls to \functionMaximizationOracleStrongCover{p}.
    %Let $X = \functionMaximizationOracleStrongCover{p}\left((G_0,w_0), S_0, T_0, y_0\right)$ be a $(p, H, w)$-tight set for some choices of inputs $(G_0, w_0), S_0, T_0,$ and $y_0$. Then, we can find a $(p, H, w)$-tight set $Y \subseteq V$ such that $X \subseteq Y$ (if such a set $Y$ exists) by guessing a vertex $u \in Y - X$ (by enumerating over all vertices $u \in V - X$) and computing $\functionMaximizationOracleStrongCover{p}\left((H,w), S\cup \{u\}, T, y\right)$. Iteratively applying this procedure will result in a maximal $(p, H, w)$-tight set. This requires at most $O(|V|)$ calls to \functionMaximizationOracleStrongCover{p}.
    \end{proof}

The next lemma says that \Cref{alg:WeakToStrongCover} can be implemented so that a single recursive call (excluding the time of recursion) runs in strongly polynomial time and makes $\poly(|V|)$ calls to \functionMaximizationOracleStrongCover{p}.
  \begin{lemma}\label{lem:WeakToStrongCover:single-recursive-call-strongly-polytime}
      Let $p:2^V\rightarrow\Z$ be a symmetric skew-supermodular function and let $\left(H = (V, E), w:E \rightarrow\Z_+\right)$ be a hypergraph satisfying conditions \ref{cond:WeakToStrong:(a)} and \ref{cond:WeakToStrong:(b)} such that $b_{(H,w)}(X) \geq p(X)$ for every $X \subseteq V$. Then, 
 \Cref{alg:WeakToStrongCover} can be implemented such that a recursive call of the algorithm on the input $((H,w), p)$ runs in strongly polynomial time and makes $\poly(|V|)$ calls to \functionMaximizationOracleStrongCover{p}.
 %, where $p:2^V\rightarrow \Z$ is a symmetric skew-supermodular function and $(H, w)$ is a hypergraph such that $b_{(H,w)}(Z) \geq p(Z)$ for every $Z \subseteq V$, $b_{(H,w)}(u)>0$ for every $u \in V$, and $\maximalTightSetContainingHyperedge{p}(\Tilde{e}) \not = \emptyset$ for every $\Tilde{e} \in E$.
  \end{lemma}  
  \begin{proof}
  It suffices to show that  \Cref{alg:WeakToStrongCover:(4)}, \Cref{alg:WeakToStrongCover:(5)} and \Cref{alg:WeakToStrongCover:(7)} of \Cref{alg:WeakToStrongCover} can be implemented to run in strongly polynomial time.
  We recall that step \Cref{alg:WeakToStrongCover:(4)} is to compute two hyperedges $e, f \in E$ such that $\maximalTightSetContainingHyperedge{p}(e) \not = \maximalTightSetContainingHyperedge{p}(f)$. We now show how to implement this computation in strongly polynomial time.  By \Cref{lem:weak-to-strong:uncrossing:maximal-tight-set-family-disjoint}, the family $\maximalTightSetFamily{p, H, w}$ is a disjoint family. Thus, $|$\maximalTightSetFamily{p, H, w}$| = O(|V|)$. Furthermore,  by \Cref{lem:WeakToStrong::maximal-tight-set-family-computed-in-strongly-polytime}, we can explicitly compute the family $\maximalTightSetFamily{p, H, w}$ in strongly polynomial time. We note that $\maximalTightSetContainingHyperedge{p}(e) \in \maximalTightSetFamily{p, H, w}$ for every $e \in E$.
      Moreover, by \Cref{cor:weak-to-strong:uncrossing:exist-two-hyperedges-contained-in-disjoint-maximal-tight-sets}, there exist hyperedges $e, f \in E$ such that $\maximalTightSetContainingHyperedge{p}(e) \cap \maximalTightSetContainingHyperedge{p}(f) = \emptyset$. Thus, we can find the required hyperedges for step (4) of \Cref{alg:WeakToStrongCover} by finding a pair of sets $X, Y \in \maximalTightSetFamily{p, H, w}$ that each contain a hyperedge, and define these hyperedges to be $e$ and $f$ respectively.

      We now show how to implement \Cref{alg:WeakToStrongCover:(5)} of \Cref{alg:WeakToStrongCover} in strongly polynomial time. In particular, we show how to compute $\beta^M$ in strongly polynomial time.
We have the following:
\begin{align*}
    \beta^M& = \min\left\{b_{(H,w)}(X) - p(X) : X \subseteq V \text{ and } e \cap X, f \cap X \not = \emptyset\right\} &\\
    & = \min_{u \in e, v \in f}\left\{\min\left\{b_{(H,w)}(X) - p(X) : \{u,v\} \subseteq X \subseteq V\right\}\right\}&\\
    &=: \min_{u \in e, v \in f}\left\{\beta^M_{uv}\right\},&
\end{align*}
where the quantity $\beta^M_{uv}$ is defined by the third equation. Then, $\beta^M$ can be computed in strongly polynomial time because  for every $u \in e$ and $v \in f$, the value $\beta^M_{uv}$ can be computed using a single oracle call to \functionMaximizationOracleStrongCover{p}. Finally, we observe that step \Cref{alg:WeakToStrongCover:(7)} of the algorithm can be computed in strongly polynomial time using a single call to $\functionMaximizationOracleStrongCover{p}$.
  \end{proof}

    We now show the main result of this section which says that \Cref{alg:WeakToStrongCover} can be implemented to run in strongly polynomial time.

\begin{lemma}\label{lem:WeakToStrongCover:strongly-polynomial-runtime:main}
    Let $p:2^V\rightarrow\Z$ be a symmetric skew-supermodular function and let $\left(H = (V, E), w:E \rightarrow\Z_+\right)$ be a hypergraph satisfying conditions \ref{cond:WeakToStrong:(a)} and \ref{cond:WeakToStrong:(b)} such that $b_{(H,w)}(X) \geq p(X)$ for every $X \subseteq V$. Then, \Cref{alg:WeakToStrongCover} can be implemented such that an execution of the algorithm on the input $((H,w), p)$ terminates in strongly polynomial time and makes only $\poly(|V|)$ calls to \functionMaximizationOracleStrongCover{p}.
    %, where $p:2^V\rightarrow \Z$ is a symmetric skew-supermodular function and $(H, w)$ is a hypergraph such that $b_{(H,w)}(Z) \geq p(Z)$ for every $Z \subseteq V$, $b_{(H,w)}(u)>0$ for every $u \in V$, and $\maximalTightSetContainingHyperedge{p}(\Tilde{e}) \not = \emptyset$ for every $\Tilde{e} \in E$.
\end{lemma}
\begin{proof}
 By \Cref{lem:WeakToStrongCover:NumberOfHyperedges} the execution of \Cref{alg:WeakToStrongCover} witnesses at most $\ell = O(|E| + |V|)$ recursive calls. Let $i \in [\ell]$ be a recursive call and let the hypergraph $(H_{i}, w_i)$ and function $p_i$ denote the inputs to the $i^{th}$ recursive call. Then, by \Cref{lem:WeakToStrong:p'-oracle-from-p-oracle}, \functionMaximizationOracleStrongCover{p_i} can be used to implement \functionMaximizationOracleStrongCover{p_{i+1}} in strongly polynomial time. By induction on $i$, we have that \functionMaximizationOracleStrongCover{p_i} can be implemented using \functionMaximizationOracleStrongCover{p} in strongly polynomial time. The claim then follows by \Cref{lem:WeakToStrongCover:single-recursive-call-strongly-polytime}.
\end{proof}
\fi


